%%%%%%%%%%%%%%%%%%%%%%%%%%%%%%%%%%%%%%%%%%%%%%%%%%%%%%%%%%%%%%%%%%%%%%
% njuthesis 示例模板 v1.4.3 2025-05-21
% https://github.com/nju-lug/NJUThesis
%
% 贡献者
% Yu XIONG @atxy-blip   Yichen ZHAO @FengChendian
% Song GAO @myandeg     Chang MA @glatavento
% Yilun SUN @HermitSun  Yinfeng LIN @linyinfeng
% Yukai CHOU @Muzimuzhi Liu Dongmiao@liudongmiao
% Wenrui HUANG @rijuyuezhu
%
% 许可证
% LaTeX Project Public License（版本 1.3c 或更高）
%%%%%%%%%%%%%%%%%%%%%%%%%%%%%%%%%%%%%%%%%%%%%%%%%%%%%%%%%%%%%%%%%%%%%%

%---------------------------------------------------------------------
% 一些提升使用体验的小技巧：
%   1. 请务必使用 UTF-8 编码编写和保存本文档
%   2. 请务必使用 XeLaTeX 或 LuaLaTeX 引擎进行编译
%   3. 不保证接口稳定，写作前一定要留意版本号
%   4. 以百分号(%)开头的内容为注释，可以随意删除
%---------------------------------------------------------------------

%---------------------------------------------------------------------
% 请先阅读使用手册：
% http://mirrors.ctan.org/macros/unicodetex/latex/njuthesis/njuthesis.pdf
%---------------------------------------------------------------------

\documentclass[
    % 模板选项(注意右端逗号)：
    %
    type = bachelor,
    % degree = academic|professional,        % 学位类型，默认为学术型
    %
    % nl-cover,   % 是否需要国家图书馆封面，默认关闭
    decl-page,  % 是否需要诚信承诺书或原创性声明，默认关闭
    %
    %   页面模式，详见手册说明
    % draft,                  % 开启草稿模式
    % anonymous,              % 开启盲审模式
    % minimal,                % 开启最小化模式
    %
    %   单双面模式，默认为适合印刷的双面模式
    % oneside,                % 单面模式，无空白页
    twoside,                % 双面模式，每一章从奇数页开始
    %
    %   字体设置，不填写则自动调用系统预装字体，详见手册
    % fontset = win|mac|macoffice|fandol|none,
  ]{njuthesis}

% 模板选项设置，包括个人信息、外观样式等
% 较为冗长且一般不需要反复修改，我们把它放在单独的文件里
\input{njuthesis-setup.def}


% 自行载入所需宏包
% \usepackage{subcaption} % 嵌套小幅图像，比 subfig 和 subfigure 更新更好
% \usepackage{siunitx} % 标准单位符号
% \usepackage{physics} % 物理百宝箱
% \usepackage[version=4]{mhchem} % 绘制分子式
\usepackage{listings} % 展示代码
\usepackage{tikz}
\usetikzlibrary{positioning,arrows.meta,shapes.geometric,shapes.symbols}
\usepackage{pgfplots}
\pgfplotsset{compat=1.18}
\usepackage{algorithm} % 浮动体
\usepackage{algpseudocode} % 伪代码（algorithmicx）
\floatname{algorithm}{算法}
\renewcommand{\algorithmicrequire}{输入：}
\renewcommand{\algorithmicensure}{输出：}
\makeatletter
\@addtoreset{equation}{chapter}
\@addtoreset{figure}{chapter}
\@addtoreset{table}{chapter}
\@addtoreset{algorithm}{chapter}
\makeatother
\renewcommand{\theequation}{\thechapter.\arabic{equation}}
\renewcommand{\thefigure}{\thechapter.\arabic{figure}}
\renewcommand{\thetable}{\thechapter.\arabic{table}}
\renewcommand{\thealgorithm}{\thechapter.\arabic{algorithm}}

% 在导言区随意定制所需命令
% \DeclareMathOperator{\spn}{span}

% 开始编写论文
\begin{document}

%---------------------------------------------------------------------
%	封面、摘要、前言和目录
%---------------------------------------------------------------------

% 生成封面页
\maketitle

% 模板默认使用 \flushbottom，即底部平齐
% 效果更好，但可能出现 underfull \vbox 信息
% 以下命令用于抑制这些信息
\raggedbottom

\begin{abstract}

在插入和删除时 k-kick 哈希通过允许有限次元素搬移，在保持查询的常数时间复杂度的同时降低了每键额外空间开销。本文先实现该结构的简化原型，再在其上加入由热 tail cubby 和 SQLite3 组成的冷热分层存储方案。基础原型保留 facility 分区、tail cubby、local query router 和商化压缩组件。该实现用于验证有限踢出、局部路由和键恢复能否在同一系统中配合运行。多层级 cubby 需要维护合并、拆分和删除回填状态，直接持久化每一层结构并不方便。因此，冷热分层方案将近期插入的数据暂存在内存 tail cubby；tail cubby 写满后，再以批量事务写入 SQLite3，冷数据按键值记录管理。实验设置多层级内存结构、纯 SQLite3 结构和冷热分层结构三类对照。各组使用相同键集合、访问序列和统计口径，比较操作延迟与平均每键空间占用。结果显示，在写后读和近期访问负载下，内存热层缩短了部分查询路径；删除近期键时，也减少了多层级回填；与纯 SQLite3 相比，冷热分层方案的查询延迟、删除延迟和按本文口径统计的每键空间占用更低，但部分规模下写入延迟并不占优。

\end{abstract}

\begin{abstract*}

By allowing a bounded number of element relocations during insertion and deletion operations, k-kick hashing reduces per-key space overhead while preserving constant-time lookup performance. This thesis first implements a simplified prototype of that structure, then adds a hot-cold storage scheme built from an in-memory tail cubby and SQLite3. The base prototype keeps facility partitioning, the tail cubby, the local query router, and quotient compression. These components are used to test whether bounded kicking, local routing, and key recovery can work together in one system. Direct persistence of the original multi-tier cubby structure is challenging, because merging, splitting, and deletion backfilling all depend on current tier state. The hot-cold design therefore keeps recently inserted keys in the in-memory tail cubby; when the tail becomes full, it flushes valid keys to SQLite3 in a batch transaction, and cold keys are managed as database records. The experiments compare the multi-tier in-memory structure, a pure SQLite3 structure, and the hot-cold tiered structure under the same key sets, access sequences, and metric definitions. Under write-after-read and recent-access workloads, the memory hot layer shortens part of the query path, and deleting recent keys triggers less multi-tier backfilling. Compared with pure SQLite3, the hot-cold scheme has lower query latency, deletion latency, and per-key space usage under the measurement method used here, but its insertion latency is not lower at all scales.

\end{abstract*}

% 生成目录
\tableofcontents
% 生成图片清单
% \listoffigures
% 生成表格清单
% \listoftables

%---------------------------------------------------------------------
%	正文部分
%---------------------------------------------------------------------
\mainmatter

% 建议将论文内容拆分为多个文件，按章节组织在 chapters/ 下 
\chapter{引言}

\section{实验背景及意义}

哈希表通过哈希函数把关键字映射到数组位置\cite{bender2024sosa}，常用于数据库索引、键值存储、编译器符号表和运行时系统\cite{chang2006bigtable,dong2021rocksdb}。在这类场景中，查询路径通常较短，工程实现也相对成熟。但当数据规模扩大后，哈希表为维持查询和更新所保存的辅助信息会逐渐占据不可忽略的空间。系统除存放键或键的部分信息外，还要记录冲突处理、删除标记、元素重定位和键恢复所需的元数据。面向键值存储的哈希索引优化也是国内已有工作的关注方向之一，相关研究从键摘要压缩、热点感知调度等角度改进空间利用和探查路径\cite{wang2024hotaware}。因此，评价哈希表时需要同时考察平均查询时间和每个键背后的元数据成本。

相较于链地址法，开放寻址哈希表把元素直接放在连续数组中，查询时不需要沿指针跳转，也更容易利用缓存局部性。它的代价也比较明确：负载因子升高后，冲突和探测长度会增加，工程实现往往需要预留空槽来降低冲突概率。删除操作还会带来额外限制。删除后的空槽如果直接释放，可能破坏后续查询路径；若采用回填或重定位，又必须维护被移动元素的状态。因此本文关注的核心问题是：在查询、插入和删除时间可控的前提下，尽量降低元素定位、更新和键恢复所需的辅助存储，并观察相关理论结构落到原型系统后产生的维护成本。

Bender 等人的论文\cite{bender2022otsh} 从理论层面给出了更详细的描述：若插入和删除允许承担 \(O(k)\) 次元素交换代价，即一次更新过程中最多搬移 \(O(k)\) 个元素，则每个元素的额外空间可以从 \(O(\log\log n)\) 比特降至 \(O(\log^{(k)}n)\) 比特，同时查询时间仍保持 \(O(1)\)，这里 \(\log^{(k)}n\) 表示对 \(n\) 连续取 \(k\) 次对数。该结果表明，开放寻址哈希表的空间开销并不只由负载率决定，还与更新过程中的元素搬移次数以及从表中恢复原始键所需的编码信息有关。本文基于这一结构完成原型实现与实验分析，重点观察上述理论权衡在工程实现中会转化为哪些具体开销。

\section{国内外研究现状}

\subsection{经典冲突处理}

常见哈希表实现大体上能分成链地址法和开放寻址法\cite{bender2024sosa}两类，链地址法一般借助链表或桶来处理冲突，写法比较直接，插入删除也相对好理解，可每个元素常要额外存指针，内存怎么分配也会打乱连续访问的效果；开放寻址法则把元素放进数组槽里，访问时数据更集中，但也更容易受负载率、探针顺序以及删除策略的影响，线性探测和它的几种改法通过连续探查换取更好的缓存表现\cite{pagh2004linear,thorup2012tabulation}，Robin Hood 哈希这类方法主要想让不同元素的查找距离尽量接近\cite{conway2018}，这些做法一到删除多或装填偏高的情形就麻烦，用墓碑占位会把查询拖长，若马上回填又得保证被挪走的元素仍能按原路径查到。

\subsection{踢出式与完美哈希}

布谷鸟哈希\cite{pagh2004cuckoo,lehmann2023sichash}会给每个元素准备多个候选槽，冲突时通过踢出把查找变短，但负载升高后仍要处理踢出链和整体重哈希这类问题；动态完美哈希\cite{dietzfelbinger1994}在最坏情况下也能把查询时间压在常数级，给了经典的上界结果，最小完美哈希\cite{lehmann2026mphsurvey,lehmann2024shockhash,hermann2024phobic,pibiri2021pthash,esposito2020recsplit}主要服务静态集合，查找保证更强，更适合关键词不再变、查多改少的场合，其后又有工作去压建表开销、压缩编码、缩短找种子的过程\cite{lehmann2025combined}，这类结构在不动态变化的数据上省空间，可一旦放进需要频繁增删的键值存储，还得另加办法才能做增量更新，重构失败时也要另行处理。

\subsection{紧凑存储与空间--时间权衡}

只看负载率还说不清紧凑哈希表到底费多少空间，因为系统还得留着辅助信息来支撑查询、更新和把键找回来，Bloom 过滤器\cite{bloom1970,broder2003bloom}这类近似判断结构，加上动态检索和过滤器下界方面的工作\cite{kuszmaul2024filters}，都说明紧凑结构一旦要支持更新，状态一变就要多占信息；紧凑检索结构\cite{pagh2008retrieval}则演示了怎样用较小空间存和键有关的值，Bender 等人提出的时间/空间权衡框架\cite{bender2022otsh}把元素放置写成带删除的球槽分配问题\cite{aamand2021dynamic,bansal2022deletions}，再用 k-kick、local query router、mini-array 和商压缩拼出一条可以调节的权衡曲线，本文的基础原型主要就是沿这一方向展开的。

\subsection{工程实现与国内研究}

打包压缩哈希表\cite{kurpicz2023pachash}、IcebergHT\cite{pandey2023iceberght,bender2024iceberg}等方案会在真实机器上测实现效果，除了吞吐和装填率，还要把缓存怎么用、重构要花多少代价这些工程因素算进去；国内也有相关工作，比如持久化键值库上的哈希索引优化\cite{wang2024hotaware}、GPU 上的动态哈希表索引\cite{xiong2022starfish}，分别从热点怎么分布、键摘要怎么压、显存怎么分页管理等方向，试着把索引占用的空间降下来。

\section{本文主要工作}

本文在 Bender 等人的论文\cite{bender2022otsh}理论框架下实现 facility--cubby--slot 分层结构的简化原型，完整理论结构中的部分细节则按原型实验需要进行了简化。facility 负责表级分区，cubby 负责局部存储和局部元数据存储，slot 存放商压缩后的高位数据。mini-array 用于维护变长元数据；cubby 内的 k-kick 完成有界交换插入；local query router 将查询从探测序列扫描转化为局部映射访问；mini-array \(B\) 保存键恢复所需的差分信息。该原型保留地址拆分、有限踢出、局部路由和键恢复四个环节。通过实验，可以分别测量查询路径、元素搬移和层级重构带来的开销。

原始多层级 cubby 结构在落地时需要维护层级合并、拆分和删除回填。在完成基础原型后，本文进一步补充冷热分层键存储方案：tail cubby 作为近期数据入口，写满后把较冷的键批量下沉到 SQLite3，这里较冷的键指的是其访问频次相对较低。这个改造借鉴了日志结构合并树先写内存、再批量下沉的思路\cite{oneil1996lsm,dong2021rocksdb}，但本文实现不再完整地保留原结构中多层级 cubby 结构，只保留 tail cubby。第四章用于检验多层级原型结构是否能按规则运行，第五章则单独讨论冷热分层改造与两个对照结构之间的差异。

围绕上述实现，实验部分关注四个问题：(1) 理论模块如何落到具体数据结构，并在插入、查询、删除过程中保持状态一致；(2) mini-array、k-kick 与 local query router 实现并组成完整的哈希系统后，插入、查询与删除各阶段的延迟处于何种量级；(3) 在不同的 \(n\)、\(K\)、\(k\) 分组实验下，k-kick 踢出次数、cubby 拆分与合并次数等指标如何变化，单次搬移次数是否仍受 \(k\) 控制；(4) 冷热分层改造后，内存 tail cubby 与 SQLite3 冷数据层在写入、查询、删除和空间占用效果如何。

\section{本文组织结构}

本章介绍实验背景及意义、国内外哈希表研究的现状以及本文的工作。第\ref{ch:theo_base}章给出相关理论基础以及冷热分层存储思想；第\ref{ch:system}章说明 facility、mini-array、cubby、k-kick、local query router 与键恢复过程的实现细节和冷热分层存储的设计方案。第\ref{ch:exp}章测量多层级原型结构在不同参数下的延迟、踢出和层级重构行为；第\ref{ch:hotcold}章讨论冷热分层改造在写后读和近期访问情况下，插入、删除、查询和空间占用的表现。第\ref{ch:concl}章总结本文工作，并阐述未来优化方向。

\chapter{问题模型与理论基础}
\label{ch:theo_base}

\section{探测复杂度与球槽分配模型}

开放寻址哈希表可被抽象为在线状态下的球槽分配问题\cite{aamand2021dynamic,bansal2022deletions,mitzenmacher2001}：表中有固定数量的槽位，每个槽位至多容纳一个元素；系统需要支持动态插入与删除，且任一时刻元素个数不超过槽位数。对任意键 \(x\)，哈希过程会为其生成一条预先确定的随机探测序列
\[
  h_1(x), h_2(x), \ldots
\]
元素只会被安置于该序列给出的候选槽位之一。

当某元素最终被置于其探测序列的第 \(i\) 个候选位置时，探测复杂度（查询时为记录“元素落在第 \(i\) 个候选位置”所需的比特数，记为 \(\Theta(1+\log i)\)，与辅助元数据空间直接相关）定义为 \(\Theta(1+\log i)\)。该量可理解为查询阶段为记录“元素采用了第几个候选位置”而必须保存的比特数目，因而与辅助元数据的空间开销直接相关。在该模型下，用平均探测复杂度刻画整体附加空间，用交换成本刻画单次更新中被搬移元素的个数；最优哈希表理论所关注的核心问题，是在允许有限交换的前提下，如何尽可能降低平均探测复杂度\cite{bender2022otsh,kuszmaul2024filters,pagh2004linear}。一般而言，若表内 \(n\) 个元素分别落在其探测序列的第 \(i_1,i_2,\ldots,i_n\) 个候选位置，则平均探测复杂度为
\[
  \frac{1}{n}\sum_{t=1}^{n}\Theta(1+\log i_t),
\]
而单次插入或删除的交换成本为踢出链中被搬移元素个数的上界。Bender 等人\cite{bender2022otsh} 通过调节 k-kick 深度参数 \(k\)，使上述两项复杂度分别控制在 \(O(\log^{(k)}n)\) 与 \(O(k)\) 量级，形成可调的时间/空间权衡曲线。他们提出：若查询必须在常数时间内完成，且插入、删除允许承担 \(O(k)\) 次元素交换，则每个元素所需的额外空间可由 \(O(\log\log n)\) 比特降至 \(O(\log^{(k)}n)\) 比特。该结果说明，开放寻址哈希表的空间效率不仅取决于空位占比，还与元素搬移、局部引导信息及键恢复编码密切相关。后文将给出与之对应的地址拆分、分层存储与 Local Query Router 设计，使探测序列在插入阶段完整定义、同时满足查询阶段不必线性扫描。

\section{系统核心设定：全域与哈希基础}

\subsection{全局参数}

设键空间全域为 \(U\)，且 \(|U|=\mathrm{poly}(n)\)。当前表内元素规模提示为 \(n\)；表容量 \(N\) 取为 2 的幂，使实际容量落在区间 \([N,2N]\) 内。每个 Facility 的基准规模为
\[
  K=\mathrm{polylog}\,n,
\]
通常取 \(K=\Theta(\log^c n)\) 且 \(c\in[2,4]\)，以便在 Facility 数量 \(N/K\) 与局部结构规模之间取得平衡\cite{bender2022otsh,bender2024sosa}。

\subsection{可逆置换与全局路由哈希}

对任意键 \(x\)，先经随机可逆置换函数 \(\pi\) 映射为 \(\pi(x)\)。该类置换可由 Feistel 等构造从伪随机函数生成\cite{luby1988,carter1979,thorup2012tabulation}。定义 \(x[a, b]\) 表示 \(x\) 的第 \(a\) 位到第 \(b\) 位组成的子串（从低位到高位），则全局路由哈希 \(g^\ast(x)\) 定义为 \(\pi(x)\) 的低 \(\log N\) 位：
\[
  g^\ast(x)=\pi(x)[1,\log N].
\]
置换值在 \(\log N\) 位以上的部分存入槽位载荷，记为
\[
  \mathrm{storage}=\pi(x)[\log N+1,\log U],
\]
即商压缩意义下的 payload；该部分不再参与 Facility 与 Cubby 的路由，仅在命中校验时与局部元数据联立以恢复完整的 \(\pi(x)\)。

从比特布局上看，\(\pi(x)\) 可示意划分为“高位 payload + 低位 \(g^\ast(x)\)”；而 \(g^\ast(x)\) 自身又可自高向低继续划分为 Facility 编号段、Cubby 中间段与 Local Query Router 段，分别用于表级分区、Cubby 内桶化与 \(A[i]\) 索引。图~\ref{fig:gx-layout} 概括了这一层次。

\begin{figure}[htbp]
  \centering
  \setlength{\tabcolsep}{4pt}
  \small
  \begin{tabular}{|c|c|c|} 
    \hline
    \multicolumn{1}{|c|}{$\pi(x)$ 高位：payload / storage} &
    \multicolumn{2}{c|}{$g^\ast(x)=\pi(x)[1,\log N]$} \\
    \hline
    $\pi(x)[\log N+1,\log U]$ & Facility 位 & Cubby 中间位 + $g_I(x)$ \\
    \hline
    存入槽位 storage & $r=g^\ast[\log K+1,\log N]$ & $g^\ast[1,\log K]$ \\
    \hline
  \end{tabular}
  \caption{$\pi(x)$ 与 $g^\ast(x)$ 的位段划分}
  \label{fig:gx-layout}
\end{figure}

\subsection{Facility 编号与 Cubby 内局部地址}

在 \(g^\ast(x)\) 已确定的前提下，Facility 编号取
\[
  r=g^\ast(x)[\log K+1,\log N]
  \;=\;\left\lfloor \frac{g^\ast(x)}{K}\right\rfloor .
\]
对容量为 \(|I|\) 的 Cubby，其内部 Local Query Router 段（与 k-kick 偏好序列无关）定义为
\[
  g_I(x)=g^\ast(x)[1,\log |I|],
\]
即 \(g^\ast(x)\) 的最低 \(\log|I|\) 位，取值范围为 \([0,|I|-1]\)。所有满足 \(g_I(x)=i\) 的键由该 Cubby 内第 \(i\) 号 Local Query Router \(A[i]\) 管理；\(g^\ast(x)\) 中 \([{\log|I|+1},{\log K}]\) 的中间位则与 Facility 编号 \(r\) 及 payload 共同用于恢复完整 \(g^\ast(x)\)，进而经由 \(\pi^{-1}\) 得到原始键 \(x\)\cite{bender2022otsh,pagh2008retrieval}。

\section{核心组件层级关系}

在逻辑结构上，整张表可视为一个大规模数组，均等划分为 \((N/K)\) 个 Facility；每个 Facility 管理规模为 \(K\) 的局部范围，并在其内部维护多层级、不同容量的 Cubby 集合。Cubby 是实际存放元素的局部块；每个 Cubby 内部又包含以下部分：
\begin{itemize}
  \item storage 数组：长度为 \(|I|\) 的槽位数组，存放商压缩 payload 即 \(\pi(x)[\log N+1,\log U]\)；
  \item k-kick 逻辑：将整个 cubby 逻辑上划分为多层级 bin，完成插入、踢出与被踢元素的重插，不要求额外的物理树结构；
  \item MiniArray \(A\)：长度为 \(|I|\) 的 Local Query Router 数组，\(A[i]\) 维护所有 \(g_I(x)=i\) 的键到探测序列下标的映射；
  \item MiniArray \(M\)：记录各 k-kick 深度 bin 的空闲槽位信息，支持在常数时间内判定 bin 内是否存在空位、并定位空位位置；
  \item MiniArray \(B\)：与 storage 槽位一一对应，保存 \(\delta=g_I(x)-i\) 及恢复 \(g^\ast(x)\) 所需的中间位等信息。
\end{itemize}

Facility 层另维护一棵基于 B 树、深度为 \(O(1)\) 的 MiniArray，用于在变长元数据与 Facility 级索引上实现均摊常数时间的访问、插入、删除与更新\cite{bender2022otsh,bayer1972}。MiniArray 中每个内部节点最多有多少个子节点称为 fanout，取 \(\Theta(\log n/\log\log n)\)，内部指针仅需 \(\Theta(\log\log n)\) 比特，从而到达整体 \(O(\log^{(k)}n)\) 的空间目标\cite{pagh2001,jacobson1989}。

为在动态扩缩容下维持高装填因子，每个 Facility 还引入 \emph{tail Cubby} 机制：初始时仅含一个 1-tiered Cubby 作为 tail；此后所有新插入默认进入 tail，tail 满后转为普通 1-tiered Cubby 并新建 tail。不同 \(j\)-tiered Cubby 的目标个数与容量分别为
\[
  t_j=\Theta\!\left(\frac{(\log^{(j+1)}n)^2}{(\log^{(j)}n)^2}\right),\qquad
  r_j=\frac{K}{(\log^{(j)}n)^2};
\]
当某层 Cubby 数量偏离 \(t_j\) 时，通过向上合并或向下拆分并在新 Cubby 内重插全部元素来恢复分布不变量。删除非 tail 元素时，从 tail 随机抽取一个元素回填至被删槽位；tail 为空时从 1-tiered 层提升一个 Cubby 作为新 tail，并可能触发拆分。tier 调度、tail 不变量与 k-kick 踢出规则共同构成第三章的算法主体\cite{bender2022otsh,bender2024sosa,conway2018,pandey2023iceberght}。

\begin{figure}[htbp]
  \centering
  \begin{tikzpicture}[
    node distance=0.65cm,
    box/.style={draw, rounded corners=2pt, minimum height=0.7cm, align=center, font=\small, inner sep=3pt},
    top/.style={box},
    child/.style={box, minimum width=2.55cm, minimum height=0.82cm},
    arr/.style={-{Stealth[length=2mm]}, thick}
  ]
    \node[top, minimum width=3.8cm] (ht) {哈希表};
    \node[top, minimum width=6cm, below=of ht] (fac) {Facility（$N/K$ 个表级分区）};
    \node[top, minimum width=6cm, below=of fac] (cub) {Cubby（多层级 + tail）};
    \node[child, anchor=north] (kick) at ([xshift=-4.5cm, yshift=-1.55cm]cub.south) {k-kick\\逻辑 bin};
    \node[child, anchor=north] (router) at ([xshift=-1.5cm, yshift=-1.55cm]cub.south) {Local Query Router\\（$A[i]$ 前缀树）};
    \node[child, anchor=north] (M) at ([xshift=1.5cm, yshift=-1.55cm]cub.south) {MiniArray $M$\\（空闲位图）};
    \node[child, anchor=north] (B) at ([xshift=4.5cm, yshift=-1.55cm]cub.south) {MiniArray $B$\\（键恢复差分）};
    \draw[arr] (ht) -- (fac);
    \draw[arr] (fac) -- (cub);
    \coordinate (bus) at ([yshift=-0.55cm]cub.south);
    \draw[arr] (cub.south) -- (bus);
    \draw[arr] (bus) -| (kick.north);
    \draw[arr] (bus) -| (router.north);
    \draw[arr] (bus) -| (M.north);
    \draw[arr] (bus) -| (B.north);
  \end{tikzpicture}
  \caption{Facility--Cubby--Slot 分层架构}
  \label{fig:sys-arch}
\end{figure}

\section{商压缩的信息论思想}

如果在槽位中完整保存原始键，则部分信息与寻址过程重复。商压缩借助双射函数 \(\pi\)，映射 \(x\) 为 \(\pi(x)\)，低 \(\log N\) 位 \(\pi(x)[1,\log N]=g^\ast(x)\) 的各位段交由地址上下文（Facility 编号、Cubby 容量、槽位下标）与局部元数据分担，槽位仅保存 payload 即 \(\pi(x)[\log N+1,\log U]\) 与无法由上下文恢复的差分信息。

具体地，当 storage\([i]\) 存放键 \(x\) 对应信息时，MiniArray \(B[i]\) 保存
\[
  \delta=g_I(x)-i,
\]
以及 \(g^\ast(x)[\log|I|+1,\log K]\) 的比特位。配合 Facility 编号 \(r=g^\ast(x)[\log K+1,\log N]\) ，即可重建
\[
  g^\ast(x)[1,\log N],
\]
再加上 payload 得到 \(\pi(x)\)，经 \(\pi^{-1}\) 恢复 \(x\)。该编码方式与紧凑检索结构\cite{pagh2008retrieval,lehmann2026mphsurvey}、打包式哈希表\cite{kurpicz2023pachash,hermann2024phobic} 所强调的“按条目实际长度计费”相同；但在 k-kick 动态搬移与 tier 重建下，\(B\)、\(A\) 与 \(M\) 必须同步更新，对 MiniArray 与 Local Query Router 的动态维护提出了比静态紧凑哈希更严格的要求\cite{bender2022otsh,kuszmaul2024filters}。

\section{k-kick、MiniArray 与 Local Query Router}

\subsection{k-kick：逻辑分层与交换上界}

在 Cubby 内部，k-kick 并不引入物理上的多层数组，而是将 storage 上长度为 \(|I|\) 的连续槽位按固定规则划分为 \(k+1\) 个逻辑深度：深度 0 的 bin （ bin 指 Cubby 按不同深度划分的子区间）覆盖整个 Cubby，更深层 bin 由上一层均匀切分得到，规模 \(s_i=\Theta((\log^{(i)}K)^6)\)。对键 \(x\)（此处均指 \(\pi(x)\)）定义偏好序列
\[
  g_0(x)\rightarrow g_1(x)\rightarrow\cdots\rightarrow g_k(x),
\]
其中 \(g_0(x)\) 为整个 Cubby，\(g_{i+1}(x)\) 为 \(g_i(x)\) 经均匀切分后 \(x\) 所属的子 bin。探测位置序列 \(h_1(x),h_2(x),\ldots\) 按先深后浅顺序枚举各 bin 内槽位；若元素最终落在第 \(j\) 个候选位置，则其探测复杂度为 \(\Theta(1+\log j)\)。

定义 bin 在已满且其中每个元素的 insert\_depth 均不低于该 bin 对应深度时饱和（IsSaturated）；贪心策略下 GetMaxUsableDepth(\(x\)) 返回 \(x\) 可插入的、未饱和 bin 的最大深度。插入时，系统先取随机哈希 \(s(x)\in[0,k]\)，再令
\[
  \mathrm{depth}=\min\bigl(\Call{GetMaxUsableDepth}{x},\,s(x)\bigr),
\]
在 bin \(g_{\mathrm{depth}}(x)\) 内写入；若 bin 已满，则踢出其中 insert\_depth 最小的元素并调用 ReInsert 向上安置。Bender 等人\cite{bender2022otsh} 证明：在参数满足其设定时，整条踢出链至多 \(k\) 次，单次插入的交换成本为 \(O(k)\)；配合随机深度选择，平均探测复杂度可控制在 \(O(\log^{(k)}n)\) 量级，从而与 上述描述的空间—时间权衡曲线对齐\cite{pagh2004cuckoo,lehmann2023sichash,pagh2004linear}。

\subsection{MiniArray \(M\)}

k-kick 的 IsSaturated、GetMaxUsableDepth 与 ReInsert 都需要在 bin 上判断是否存在空位并找到可用位置。MiniArray \(M\) 为各深度 bin 维护等价于 bitmap 的空闲槽位图\cite{wang2023bitmap}，使上述判定与定位在均摊意义下为常数时间，而不必扫描整个 Cubby\cite{bender2022otsh,pagh2004linear}。 MiniArray 则通过 \(O(1)\) 深度 B 树、叶内 Rank/Select 与 Packed Leaf 维护变长元数据，并实现常数时间复杂度的 Access/Insert/Delete/Update 函数\cite{bayer1972,jacobson1989,pagh2001}：树高 \(O(1)\)，内部指针 \(\Theta(\log\log n)\) 比特\cite{kurpicz2023pachash,lehmann2024shockhash}。

\subsection{Local Query Router：\(O(1)\) 查询与映射存储}

偏好序列与探测位置序列在插入阶段完整定义；为实现 $O(1)$ 查询，须额外维护键到探测下标的映射，而不遍历 $h_1(x), h_2(x), \dots$：k-kick 插入会搬移元素，其实际位置既不等于首选位置 $g_k(x)$，也不等于 $h_1(x)$。Local Query Router 因此为每个键保存映射
\[
  \pi(x)\longrightarrow j,
\]
其中 \(j\) 为探测序列下标。Cubby 内令 \(g_I(x)=g^\ast(x)[1,\log|I|]\)，所有满足 \(g_I(x)=i\) 的键由 router \(A[i]\) 管理；\(A[i]\) 实现为二叉前缀树（可进一步 Patricia 压缩\cite{morrison1968,lehmann2025combined}），对外接口为 insert(key, j)、query(key)、delete(key)。查询时先算 \(i=g_I(x)\)，在 \(A[i]\) 得 \(j\)，再访问 storage\([h_j(x)]\)，加上 \(B[h_j(x)]\) 与 payload 恢复 \(\pi(x)\)，与查询侧的 \(\pi(x)\) 进行对比；\(j\le\log K\)，每条映射仅需 \(O(\log\log n)\) 比特，查询路径为常数次 router 与 MiniArray 访问\cite{bender2022otsh,pagh2008retrieval}。

\subsection{辅助结构同步不变量}

元素在 k-kick 踢出链中每发生一次搬移，须同步更新四处状态：
\begin{itemize}
  \item storage 槽位内容与 insert\_depth；
  \item \(B[i]\) 中的 \(\delta\) 与中间位；
  \item \(A[g_I(x)]\) 中 \(\pi(x)\mapsto j\) 的映射；
  \item \(M\) 中各 bin 的空闲位图。
\end{itemize}
删除非 tail 元素时的 tail 随机回填、tier 合并/拆分后的全量重插，以及表级双表扩缩容，均遵循同一同步规则——不可以仅复制 payload 而忽略 router 与 \(B\) 的上下文依赖\cite{bender2022otsh,bender2024sosa}。第三章将给出 MiniArray、k-kick 与 Local Query Router 的全部算法伪代码及统一操作语义。

\section{本章小结}

本章建立了后文所需的理论语境：球槽分配模型给出探测复杂度与交换成本的定义\cite{knuth1973sorting,aamand2021dynamic}；全域参数与 \(g^\ast(x)\) 位段划分明确了地址语义；图~\ref{fig:sys-arch} 与 Facility--Cubby--Slot(\(storage/A/M/B\)) 层级关系给出整体骨架；商压缩说明 payload 与 \(B[i]\) 的信息分工\cite{bender2022otsh,lehmann2025combined}；同时从 k-kick 交换上界、MiniArray \(M\) 的空闲判定、Local Query Router 的 \(O(1)\) 查询及四个状态同步不变量四个角度，为第三章算法设计提供理论锚点。下一章将给出各模块的完整伪代码；第四章则报告宏基准与参数扫描下的性能测量。

\chapter{哈希表结构与算法设计}
\label{ch:system}

本章依照系统核心设定，给出分层哈希表各模块组件的完整结构与算法细节。内容覆盖 Facility 、 MiniArray、多层级 Cubby、Cubby 内 k-kick 与 Local Query Router，商压缩的实现，并在最后一节统一说明插入、查询与删除的操作语义。

\section{Facility 与 MiniArray}

\subsection{Facility 的基本属性}

每个 Facility 对应全局表中的一个局部片段，规模为 \(K=\mathrm{polylog}\,N\)。Facility 管理其内部的全部 Cubby（包括各 tier 及唯一 tail），并维护 一个 Facility 层次的 MiniArray，用于定位键对应的 Cubby 。B 树的 fanout(上文有定义，fanout 为每个内部节点的子指针个数) 记为 \(F\)，取 \(F=\Theta(\log n/\log\log n)\)；树高为 \(O(1)\)，每个内部节点的子指针仅需 \(\Theta(\log\log n)\) 比特\cite{bender2022otsh,bayer1972,pagh2001}。

\subsection{MiniArray 的逻辑结构}

MiniArray 由以下部分共同组成：
\begin{center}
  \small
  MiniArray $\rightarrow$ $O(1)$ 深度 B 树 $\rightarrow$ Packed Leaf + Bitmap + Rank/Select + Lookup Table
\end{center}

叶节点 Leaf 包含：
\begin{itemize}
  \item bitmap：用于标记叶内哪些逻辑位置已有有效元素，其大小与叶容量相同；
  \item data：PackedArray，在有效位置上紧凑、无间隙地存放变长元素。
\end{itemize}

内部节点 Internal 包含：
\begin{itemize}
  \item subtree\_size$[F]$：每个子树中有效元素总数，供 Rank/Select 定位；
  \item children$[F]$：至多 \(F\) 个子女指针，每个指针 \(\Theta(\log\log n)\) 比特。
\end{itemize}

MiniArray 对外提供的接口函数为 Access(idx)、Insert(idx, val)、Delete(idx)、Update(idx, val)，其中 idx 为逻辑下标。叶容量上限为 MAX\_LEAF\_SIZE，由 NODE\_MAX\_BITS 常数 \(c\in[1,3]\) 产生\cite{bender2022otsh,kurpicz2023pachash,lehmann2024shockhash}。

\subsection{Access：按逻辑下标读取}

Access 操作的输入是逻辑下标，搜索过程自根向下逐层向下遍历，利用各层 subtree\_size 进行前缀和判断，将输入的逻辑下标映射到对应的目标叶节点；在叶子节点的内部用 Select 在 bitmap 上定位第 idx 个有效位置，再从 data 读出元素。树高为 \(O(1)\)，叶内 Select 算法通过 Lookup Table 常数时间完成\cite{jacobson1989,pagh2008retrieval,kurpicz2023pachash}， 不需要遍历比特串，获取到物理位置之后，直接从紧凑数组中读取对应的元素，并返回。

\begin{algorithm}[htbp]
  \caption{MiniArray.Access}
  \label{alg:miniarray-access}
  \begin{algorithmic}[1]
    \State 输入逻辑下标 idx
    \State cur $\gets$ root，rem $\gets$ idx
    \While{cur 为内部节点}
      \State 找到子树 $i$ 使 $\displaystyle\sum_{t<i}\mathit{subtree\_size}[t]\le \mathit{rem}<\sum_{t\le i}\mathit{subtree\_size}[t]$
      \State rem $\gets$ rem - $\displaystyle\sum_{t<i}\mathit{subtree\_size}[t]$
      \State cur $\gets$ cur.child[i]
    \EndWhile
    \State pos $\gets$ Call{Select}(cur.bitmap, rem)
    \State 返回 cur.data[pos]
  \end{algorithmic}
\end{algorithm}

\subsection{Insert：插入与叶分裂}

插入过程中，Insert 在定位叶节点的过程中维护路径栈，为后续自底向上更新 subtree\_size 提供途径；在叶节点内用 Select 函数找到插入位置，将 bitmap 对应位置赋值为 1，并向 data 写入 val。若叶内有效元素数达到叶子容量的上限，就会调用 Split 对叶子节点进行分裂操作；插入与分裂检查完成之后，自底向上对路径栈中每个内部节点的 subtree\_size 加一， 以保证后续查询与定位的正确性。

\begin{algorithm}[htbp]
  \caption{MiniArray.Insert}
  \label{alg:miniarray-insert}
  \begin{algorithmic}[1]
    \State 输入 idx, val；初始化空路径栈 stk
    \State cur $\gets$ root，rem $\gets$ idx
    \While{cur 为内部节点}
      \State stk.push(cur)
      \State 找到子树 $i$ 使 $\displaystyle\sum_{t<i}\mathit{subtree\_size}[t]\le \mathit{rem}<\sum_{t\le i}\mathit{subtree\_size}[t]$
      \State rem $\gets$ rem - $\displaystyle\sum_{t<i}\mathit{subtree\_size}[t]$
      \State cur $\gets$ cur.child[i]
    \EndWhile
    \State leaf $\gets$ cur
    \State pos $\gets$ Call{Select}(leaf.bitmap, rem)
    \State leaf.bitmap.set(pos)；leaf.data.insert(pos, val)
    \If{leaf.size $\ge$ MAX\_LEAF\_SIZE}
      \State Call{Split}(leaf)
    \EndIf
    \For{路径栈中每个内部节点 node}
      \State node.subtree\_size[i] $\mathrel{+}=1$
    \EndFor
  \end{algorithmic}
\end{algorithm}

\subsection{Delete：删除与叶合并}

Delete 同样沿 subtree\_size 定位叶节点，用 Select 函数在 Bitmap 上找到待删元素的具体位置，清除 bitmap 对应位，并从 data 移除对应的元素。若叶内有效元素过少，低于阈值，调用 Merge 与相邻兄弟叶合并，合并可能需要多次操作才能保证 B 树恢复数据结构；最后自底向上将路径上 subtree\_size 减一。

\begin{algorithm}[htbp]
  \caption{MiniArray.Delete}
  \label{alg:miniarray-delete}
  \begin{algorithmic}[1]
    \State 输入 idx；初始化路径栈 stk
    \State 沿 subtree\_size 定位至叶 leaf(同 Insert)
    \State pos $\gets$ Call{Select}(leaf.bitmap, rem)
    \State leaf.bitmap.unset(pos)；leaf.data.remove(pos)
    \If{leaf 有效元素过少（即 leaf.size < MAX\_LEAF\_SIZE / 2）}
      \State Call{Merge}(leaf)
    \EndIf
    \For{路径栈中每个内部节点 node}
      \State node.subtree\_size[i] $\mathrel{-}=1$
    \EndFor
  \end{algorithmic}
\end{algorithm}

\subsection{Update：原地覆写}

Update 仅修改已有有效位置的元素值，不改变 bitmap 中 1 的个数，也不触发 Split/Merge。定位过程与 Access 相同，在叶内 Select 后直接 leaf.data[pos] $\gets$ val。

\begin{algorithm}[htbp]
  \caption{MiniArray.Update}
  \label{alg:miniarray-update}
  \begin{algorithmic}[1]
    \State 输入 idx, val
    \State 沿 subtree\_size 定位至叶 leaf
    \State pos $\gets$ Call{Select}(leaf.bitmap, rem)
    \State leaf.data[pos] $\gets$ val
  \end{algorithmic}
\end{algorithm}

\subsection{Split 与 Merge}

Split（叶分裂）。 新建叶节点，以 leaf.size/2 为界将后半段 bitmap 与 data 迁至新叶，原叶保留前半段；在父节点 child 数组中插入新叶指针，并更新 subtree\_size 的分裂计数。

\begin{algorithm}[htbp]
  \caption{MiniArray.Split（叶分裂）}
  \label{alg:miniarray-split}
  \begin{algorithmic}[1]
    \State 输入叶 leaf
    \State new\_leaf $\gets$新建 LeafNode；mid $\gets$ leaf.size/2
    \State 将 leaf.bitmap、leaf.data 后半段搬至 new\_leaf
    \State leaf 保留前半段
    \State 在父节点 parent.child 中插入 new\_leaf，更新 parent.subtree\_size
  \end{algorithmic}
\end{algorithm}

Merge（叶合并）。 取相邻兄弟叶，将兄弟的 bitmap 与 data 并入当前叶，从父节点删去兄弟子树并刷新 subtree\_size。

\begin{algorithm}[htbp]
  \caption{MiniArray.Merge（叶合并）}
  \label{alg:miniarray-merge}
  \begin{algorithmic}[1]
    \State 输入叶 leaf；sib $\gets$相邻兄弟叶；parent $\gets$ leaf 的父节点
    \State leaf.bitmap $\mathrel{|=}$ sib.bitmap；leaf.data.append(sib.data)
    \State 从 parent.child 移除 sib，更新 parent.subtree\_size
  \end{algorithmic}
\end{algorithm}

\subsection{Rank 与 Select}

叶子节点中 bitmap 上的 Rank/Select 通过 Lookup Table 实现，这是紧凑动态结构的标准组件\cite{jacobson1989,pagh2008retrieval,bender2022otsh,lehmann2026mphsurvey}：
\begin{itemize}
  \item Rank(bitmap, \(i\))：返回 bitmap 前 \(i\) 位中 1 的个数，等于 LookupTable 对高段与低段的查表结果之和；
  \item Select(bitmap, \(k\))：返回 bitmap 中第 \(k\) 个 1 的位置，由 LookupTable\([k]\) 直接给出。
\end{itemize}

\begin{algorithm}[htbp]
  \caption{Rank / Select（查表实现）}
  \label{alg:rank-select}
  \begin{algorithmic}[1]
    \Function{Rank}{bmp, $i$}
      \State 返回 LookupTable[bmp 高段] $+$ LookupTable[bmp 低段]
    \EndFunction
    \Function{Select}{bmp, $k$}
      \State 返回 LookupTable[$k$] \Comment{第 $k$ 个 1 的位置}
    \EndFunction
  \end{algorithmic}
\end{algorithm}

\subsection{复杂度与空间性质}

Facility 级 MiniArray 的树高为 \(O(1)\)，自根至叶的每层路由仅依赖 subtree\_size 比较，叶内 Rank/Select 经 Lookup Table 常数时间完成\cite{jacobson1989,pagh2008retrieval}。因此 Access 与 Update 为最坏情况 \(O(1)\)；Insert 与 Delete 在叶分裂/合并均摊意义下亦为 \(O(1)\)\cite{bayer1972,bender2022otsh}。空间方面，每个内部节点的 \(F\) 个子指针各需 \(\Theta(\log\log n)\) 比特，叶节点以 bitmap 标记有效位置、以 PackedArray 无间隙存放变长元素，整体辅助位数与 \(O(\log^{(k)}n)\) 的全局预算相容\cite{kurpicz2023pachash,lehmann2024shockhash}。Cubby 内的 \(A\)、\(M\)、\(B\) 三组 MiniArray 复用同一套接口语义：\(A[i]\) 存 router 映射，\(M\) 存 bin 空闲图，\(B\) 与 storage 槽位一一对应存差分编码——三者在 k-kick 搬移时须原子同步更新。

\section{多层级 Cubby 设计}

\subsection{层级属性与初始状态}

Facility 内 Cubby 按 tier 分级管理。\(j\)-tiered Cubby 的目标数量为
\[
  t_j=\Theta\!\left(\frac{(\log^{(j+1)}n)^2}{(\log^{(j)}n)^2}\right),
\]
容量（槽位数）为
\[
  r_j=\frac{K}{(\log^{(j)}n)^2}.
\]
系统初始时，每个 Facility 仅包含一个 1-tiered Cubby，且该 Cubby 即为 tail。此后所有新键插入 tail；tail 满后，旧 tail 被视为一个普通的 1-tiered Cubby 挂入 tier 集合，并新建 tail 承接后续插入，该过程可能触发 tier 向上合并\cite{bender2022otsh,bender2024sosa,bender2024iceberg,pandey2023iceberght}。

\subsection{动态调整：合并与拆分}

设当前 \(j\)-tiered Cubby 的实际个数为 \(\mathrm{cnt}_j\)。调度规则如下：
\begin{itemize}
  \item 当 \(\mathrm{cnt}_j\ge 3t_j\) 时，取出 \(t_j\) 个 \(j\)-tiered Cubby，将其内全部元素导出，合并为一个 \((j+1)\)-tiered Cubby 并重新插入；
  \item 当 \(\mathrm{cnt}_j\le t_j/2\) 且存在非空的 \((j+1)\)-tiered Cubby 时，取出一个 \((j+1)\)-tiered Cubby，拆成 \(t_j\) 个 \(j\)-tiered Cubby 并重插元素；
  \item 无论合并或拆分，均须取出 cubby 内所有元素重新插入，并同步更新对应 Cubby 中的 MiniArray \(A\)、\(M\)、\(B\) 与 storage，不能仅复制旧槽位内容（因 \(r_j\)、\(|I|\) 与商压缩上下文可能已变）\cite{bender2022otsh,conway2018,kuszmaul2024filters}。
\end{itemize}

\begin{algorithm}[htbp]
  \caption{j-tiered cubby 数量调度}
  \label{alg:tier-schedule}
  \begin{algorithmic}[1]
    \For{每个 tier 层级 \(j\)}
      \State 计算目标数量 \(t_j\)，统计实际数量 \(\mathrm{cnt}_j\)
      \If{\(\mathrm{cnt}_j\ge 3t_j\)}
        \State 取出 \(t_j\) 个 \(j\)-tiered Cubby，合并为一个 \((j+1)\)-tiered Cubby 并重新插入
      \ElsIf{\(\mathrm{cnt}_j\le t_j/2\) 且 tier \(j+1\) 非空}
        \State 取出一个 \((j+1)\)-tiered Cubby，拆成 \(t_j\) 个 \(j\)-tiered Cubby 并重新插入
      \EndIf
    \EndFor
  \end{algorithmic}
\end{algorithm}

\subsection{Cubby 内部组件}

每个容量为 \(|I|\) 的 Cubby 包含以下五部分：

storage 数组。 长度为 \(|I|\)，存放键值的商压缩载荷；槽位 \(i\) 上的 storage\([i]\) 对应 \(\pi(x)[\log N+1,\log U]\) 的 payload 部分，完整键恢复依赖 \(B[i]\) 与全局上下文。

k-kick 插入逻辑。 在 storage 上按一定的规则动态决定元素实际槽位，并维护每个占用元素的 insert\_depth。

MiniArray \(A\)。 长度为 \(|I|\) 的 Local Query Router 数组；\(A[i]\) 是一棵二叉前缀树，管理所有满足 \(g_I(x)=i\) 的键。

MiniArray \(M\)。 存储各 k-kick 深度 bin 的空闲槽位图，使插入、删除与 ReInsert 能在常数时间内判定 bin 是否含空位、定位空位，而不必扫描整个 Cubby。

MiniArray \(B\)。 与 storage 槽位一一对应。当 storage\([i]\) 存放键 \(x\) 的信息时，\(B[i]\) 保存
\[
  \delta=g_I(x)-i,\qquad g_I(x)=g^\ast(x)[\log|I|+1,\log K],
\]
以及 \(g^\ast(x)[\log|I|+1,\log K]\) 的中间位。配合 Facility 编号 \(r=g^\ast(x)[\log K+1,\log N]\) 与 payload，可完整恢复 \(\pi(x)\)\cite{bender2022otsh,pagh2008retrieval}。

\subsection{插入与删除操作}

插入。 每次插入均进入当前 Facility 的 tail Cubby。若 tail 仍有空位，则在 tail 内执行 k-kick 插入；若 tail 已满，则新建 tail，并将已满的旧 tail 视作 1-tiered Cubby 挂入 1-tier Cubby 集合，该步骤可能因 \(\mathrm{cnt}_1\ge 3t_1\) 而触发向上合并。

删除与 tail 回填。 若删除发生在 tail 之外的其他 Cubby，删除完成后从 tail 中\emph{随机}选取一个占用元素，经 k-kick 规则回填至刚空出的槽位，并更新 MiniArray \(A\)、\(M\)、\(B\) 与 storage。若 tail 在回填后为空，则从 1-tiered 层取出一个 Cubby 提升为 tail；该过程可能因 \(\mathrm{cnt}_i\le t_i\) 而触发向下拆分(甚至可能会触发多层)。上述不变量保证：除 tail 外，各 Cubby 尽量维持满载，局部波动由 tail 与 tier 调度共同吸收\cite{bender2022otsh,aamand2021dynamic,bansal2022deletions}。

\begin{algorithm}[htbp]
  \caption{tail 维护与删除回填}
  \label{alg:tail-ops}
  \begin{algorithmic}[1]
    \State 插入：写入 tail；若 tail 满则旧 tail 归入 1-tiered Cubby 并新建 tail
    \State 删除（非 tail）：释放被删槽位；从 tail 随机取元素 \(y\) 回填
    \State 对 \(y\) 执行 k-kick 写入，更新 \(A,M,B\)
    \If{tail 为空}
      \State 从 1-tiered Cubby 提升一个 Cubby 为 tail；必要时触发 tier 拆分
    \EndIf
  \end{algorithmic}
\end{algorithm}

\section{k-kick 层级插入算法}

约定：Cubby 内算法中的 \(x\) 均指 \(\pi(x)\)（置换后的值），而非原始键；符号仍记为 \(x\)，设 Cubby 容量为 \(|I|\)。

\subsection{层级 bin 定义}

k-kick 不对 storage 做物理分层，仅在逻辑上将 Cubby 划分为 \(k+1\) 个深度，\(k\) 为常数（通常 \(k\in[3,5]\)，基准取 4），定义不同层次的 bin：
\begin{itemize}
  \item 深度 0：唯一 bin 覆盖整个 Cubby，大小 \(s_0=K=\mathrm{polylog}\,n\)；
  \item 深度 1：将深度 0 的 bin 均匀切分为许多大小 \(s_1=\Theta((\log K)^6)\) 的子 bin；
  \item \( \dots \)
  \item 深度 \(i\)（\(i\ge 1\)）：bin 大小 \(s_i=\Theta((\log^{(i)}K)^6)\)，由上一层 bin 继续均匀切分得到。
\end{itemize}

定义函数 Bin(\(x,i\)) 返回键 \(x\) 在深度 \(i\) 所属 bin 的槽位区间：深度 0 时返回 \([0,|I|-1]\)；深度 \(i\ge 1\) 时在 Bin(\(x,i-1\)) 内按 \(s_i\) 均匀切分，并根据 \(x\) 的哈希比特在每一层中选定对应的一个子区间\cite{bender2022otsh,carter1979}。

\subsection{偏好序列}

对键 \(x\) 定义偏好序列
\[
  g_0(x)\rightarrow g_1(x)\rightarrow\cdots\rightarrow g_k(x).
\]
其中 \(g_0(x)=\operatorname{Bin}(x,0)\) 为整个 Cubby；\(g_{i+1}(x)\) 为 \(g_i(x)\) 经均匀切分后 \(x\) 所落子 bin。于是 \(g_{i+1}(x)\) 一定是 \(g_i(x)\) 的子 bin；沿路径自顶向下，每一层均为上一层的随机子区间。自底向上回溯 \(g_k,g_{k-1},\ldots,g_0\) 即可得到全部祖先 bin。

\begin{algorithm}[htbp]
  \caption{PreferenceSequence：偏好序列}
  \label{alg:pref-seq}
  \begin{algorithmic}[1]
    \Function{PreferenceSequence}{$x$}
      \State $g[0]\gets\Call{Bin}{x,0}$
      \For{$i=1$ to $k$}
        \State $g[i]\gets g[i-1]$ 内均匀切分后、$x$ 所属的子 bin
      \EndFor
      \State 返回 $g[0..k]$
    \EndFunction
  \end{algorithmic}
\end{algorithm}

\subsection{探测位置序列}

\(h_j(x)\) 表示 \(x\) 在 Cubby storage 中的第 \(j\) 个候选槽位。序列按先深后浅顺序构造：依次遍历 \(g_k(x),g_{k-1}(x),\ldots,g_0(x)\)，并在每个 bin 内按槽位顺序枚举全部位置。

对深度 \(i\) 的 bin \(g_i(x)\)，设
\[
  \mathrm{base}_i=s_k+s_{k-1}+\cdots+s_{i+1},
\]
则对任意 \(j\in[0,s_i-1]\)，全局探测序列中对应项为
\[
  h_{\mathrm{base}_i+j}(x)=g_i(x)\text{ 的第 }(j+1)\text{ 个槽位}.
\]
注意文档中下标公式 \(h_{(k+1-i)\cdot s_i+j}(x)\) 与 \(\mathrm{base}_i\) 表述等价：更深层的 bin 在序列中排在更前。序列总长度约为 \(\displaystyle\sum_{i=0}^{k}s_i\)；查询阶段并不线性扫描该序列，而由 Local Query Router 直接给出实际使用的下标 \(j\)\cite{bender2022otsh,pagh2004linear}。

\begin{algorithm}[htbp]
  \caption{ProbeSequence：探测位置序列（先深后浅）}
  \label{alg:probe-seq}
  \begin{algorithmic}[1]
    \Function{ProbeSequence}{$x$}
      \State $g\gets\Call{PreferenceSequence}{x}$；$\mathit{seq}\gets[\,]$
      \For{$i=k$ downto $0$}
        \State $\mathit{bin}\gets g[i]$
        \For{$\mathit{pos}$ 从 $\mathit{bin.start}$ 到 $\mathit{bin.end}$}
          \State $\mathit{seq.append(pos)}$
        \EndFor
      \EndFor
      \State 返回 $\mathit{seq}$ \Comment{$h_1(x),h_2(x),\ldots$ 依次为 $\mathit{seq}[0],\mathit{seq}[1],\ldots$}
    \EndFunction
  \end{algorithmic}
\end{algorithm}

\subsection{bin 饱和与最大可用深度}

在深度 \(d\) 下，称 bin 饱和，当且仅当 bin 已满，且其中每个元素 \(y\) 均满足 y.insert\_depth \(\ge\) d。bin 未满时不饱和；bin 已满但存在 y.insert\_depth < d 的元素时，亦不视为饱和。

GetMaxUsableDepth(\(x\)) 自 \(k\) 向 0 扫描偏好序列，返回最深的不饱和 bin 的深度；若全部饱和则返回 0(说明已满)。

\begin{algorithm}[htbp]
  \caption{IsSaturated 与 GetMaxUsableDepth}
  \label{alg:saturated}
  \begin{algorithmic}[1]
    \Function{IsSaturated}{bin, $d$}
      \If{bin 未满} \State 返回 false \EndIf
      \For{bin 中每个元素 $y$}
        \If{y.insert\_depth < d} \State 返回 false \EndIf
      \EndFor
      \State 返回 true
    \EndFunction
    \Function{GetMaxUsableDepth}{$x$}
      \State $g\gets\Call{PreferenceSequence}{x}$
      \For{$d=k$ downto $0$}
        \If{not \Call{IsSaturated}{$g[d], d$}}
          \State 返回 $d$
        \EndIf
      \EndFor
      \State 返回 $0$
    \EndFunction
  \end{algorithmic}
\end{algorithm}

\subsection{随机深度选择与 InsertKick}

插入时须记录 insert\_depth 并在各深度 bin 间均匀分布。插入前先取随机哈希 \(s(x)\in[0,k]\)，再令
\[
  \mathrm{depth}=\min(\Call{GetMaxUsableDepth}{x},\,s(x)).
\]
该式为随机深度与可用深度的 min 组合，对贪心上界形成软上限，使高深度 bin 更长时间保持未饱和。在目标 bin \(g_{\mathrm{depth}}(x)\) 中：若有空位，则写入 \(x\)、记录 x.insert\_depth = depth，并在 Local Query Router \(A[g_I(x)]\) 中登记探测下标 \(j\)；若无空位，则踢出 bin 内 insert\_depth 最小的元素 \(y\)，将 \(x\) 写入被踢槽位，再对 \(y\) 调用 ReInsert。整条踢出链至多触发 \(k\) 次，与 \(O(k)\) 交换上界一致\cite{bender2022otsh}。

\begin{algorithm}[htbp]
  \caption{InsertKick：k-kick 插入}
  \label{alg:insert-kick}
  \begin{algorithmic}[1]
    \State 输入 $x$（已为 $\pi(x)$）
    \State $\mathit{depth\_max}\gets \Call{GetMaxUsableDepth}{x}$
    \State $s_x\gets\Call{RandomHash}{x}\bmod(k+1)$
    \State $\mathit{depth}\gets\min(\mathit{depth\_max},\,s_x)$
    \State $g\gets\Call{PreferenceSequence}{x}$；$\mathit{target}\gets g[\mathit{depth}]$
    \If{$\mathit{target}$ 有空位（由 $M$ 判定）}
      \State 在任意空位写入 $x$；$x.insert\_depth\gets\mathit{depth}$
      \State 更新 $A,M,B$；返回成功
    \Else
      \State $y\gets\mathit{target}$ 中 $insert\_depth$ 最小元素；$\mathit{evict\_pos}\gets y$ 的位置
      \State 从 $\mathit{evict\_pos}$ 移除 $y$；将 $x$ 写入 $\mathit{evict\_pos}$
      \State $x.insert\_depth\gets\mathit{depth}$
      \State \Call{ReInsert}{$y$}
    \EndIf
  \end{algorithmic}
\end{algorithm}

\subsection{ReInsert：被踢元素向上放入}

插入发生踢出时，被踢元素须安置到合适深度；\textbf{ReInsert} 自 \(y.\)insert\_depth\(-1\) 向上逐层检查 Bin(\(y,d\))：若该 bin 在深度 \(d\) 下不饱和且有空位，或存在 insert\_depth 小于 \(d\) 的元素，则将 \(y\) 写入并可能继续踢出链式调用 ReInsert。深度降至 0 仍无法安置时 Cubby 已满，须新建 tail Cubby 并重试 InsertKick，与多层级 Cubby 设计衔接。MiniArray \(M\) 在此过程中的空位判定与定位均为常数时间。理论保证：在参数满足 Bender 等人\cite{bender2022otsh} 的设定时，极大概率至多 \(k\) 次 kick 后一定能安置被踢元素，如算法~\ref{alg:insert-kick} 所示。

\begin{algorithm}[htbp]
  \caption{ReInsert：被踢元素向上找空位}
  \label{alg:reinsert}
  \begin{algorithmic}[1]
    \Function{ReInsert}{$y$}
      \For{$d=y.insert\_depth-1$ downto $0$}
        \State $\mathit{bin}\gets\Call{Bin}{y,d}$
        \If{\Call{IsSaturated}{$\mathit{bin}, d$}}
          \State \textbf{continue}
        \EndIf
        \If{$\mathit{bin}$ 有空位（由 $M$ 判定）}
          \State 将 $y$ 写入空位；$y.insert\_depth\gets d$
          \State 更新 $A,M,B$；返回成功
        \ElsIf{$\mathit{bin}$ 中存在 $z.insert\_depth<d$}
          \State $z\gets\mathit{bin}$ 中 $insert\_depth$ 最小且满足 $z.insert\_depth<d$ 的元素
          \State 从 $\mathit{bin}$ 移除 $z$；将 $y$ 写入该槽位
          \State $y.insert\_depth\gets d$；更新 $A,M,B$
          \State \Call{ReInsert}{$z$}；返回
        \EndIf
      \EndFor
      \State 返回失败 \Comment{深度已降至 $0$ 仍无法安置，Cubby 已满}
      \State 新建 tail cubby；在新 tail 中对 $y$ 重试 \Call{InsertKick}{}
    \EndFunction
  \end{algorithmic}
\end{algorithm}

\subsection{复杂度与正确性}

单次 InsertKick 至多触发一条长度为 \(k\) 的踢出链：每次踢出仅调用一次 ReInsert，而被踢元素的 insert\_depth 严格下降，故链长受 \(k\) 约束，与 Bender 等人\cite{bender2022otsh} 的 \(O(k)\) 交换成本一致。

GetMaxUsableDepth 自 \(k\) 向 0 扫描，返回最深的不饱和 bin；随机深度 \(\mathrm{depth}=\min(\mathrm{depth\_max},s(x))\) 在理论上保证：若存在可写 bin，则插入不会触发深层踢出，同时使不同深度的 bin 负载趋于均衡\cite{mitzenmacher2001}。

bin 中的是否存在空位与定位空位都通过 MiniArray \(M\) 完成，避免对整个 Cubby 进行依次遍历 \(O(|I|)\) 时间复杂度的扫描；ReInsert 从 y.insert\_depth-1 向上逐层检查，每层仅访问 \(M\) 与局部 bin 状态，均摊常数时间。

元素每次写入或迁移，都需要同步更新 storage、insert\_depth、\(B\)、\(g_I(x)\) 中的 \(\pi(x)\mapsto j\) 映射对应的信息，以及 \(M\) 的空闲位图，这一同步规则是整个系统正常运行的前提，之后的操作才能正常运行。

\section{Local Query Router}

\subsection{设计目标}

偏好序列与探测位置序列 \(h_1(x),h_2(x),\ldots\) 在插入分析中完整定义了候选槽位顺序。然而 k-kick 插入会通过踢出动态挪动元素：元素实际位置一般不等于首选深度 \(g_k(x)\)，也不等于 \(h_1(x)\)。若查询仍沿 \(h_1,h_2,\ldots\) 线性扫描，时间复杂度将与探测序列长度同阶，无法达到 \(O(1)\)。因此必须为每个键保存精准映射
\[
  \pi(x)\longrightarrow j,
\]
其中 \(j\) 为探测序列中的下标，查询时直接访问下标 \(h_j(x)\) 对应的 storage 中的数据，再用 \(B\) 与 payload 三个数据恢复 \(\pi(x)\) 后与查询元素进行对比校验\cite{bender2022otsh,pagh2008retrieval}。

\subsection{映射规则}

Cubby 内定义局部首选槽位号（与 k-kick 偏好序列无关）：
\[
  g_I(x)=g^\ast(x)[1,\log|I|]\in[0,|I|-1].
\]
所有满足 \(g_I(x)=i\) 的键由第 \(i\) 号 Local Query Router \(A[i]\) 管理。查询时先计算 \(i=g_I(x)\)，再在 \(A[i]\) 中查询 \(\pi(x)\) 得到探测序列下标 \(j\)（若不存在则未命中）；随后访问 storage 中 \(h_j(x)\) 对应槽位，读取 payload 与 \(B\)，重建 \(g^\ast(x)\) 并校验是否命中。

\subsection{二叉前缀树存储与接口}

每个 \(A[i]\) 实现为一棵二叉前缀树；对外键为 \(\pi(x)\) 的比特串。接口为：
\begin{itemize}
  \item insert(key, j)：\(j\le\log K\)，仅需 \(O(\log\log n)\) 比特存储；如算法~\ref{alg:lqr-insert} 所示；
  \item query(key)：返回 \(j\) 或“不存在”；如算法~\ref{alg:lqr-query} 所示；
  \item delete(key)：删除映射。如算法~\ref{alg:lqr-delete} 所示。
\end{itemize}

树中不存完整键，仅存区分键的最短前缀；叶节点存 j\_value。

插入过程利用映射后的 \(\pi(x)\) 进行操作；下文以 \(x\) 表示 \(\pi(x)\) 的比特串。映射后的 \(x\) 的每一位（从高位到低位）依次决定在树中向左（0）还是向右（1）走。沿路径向下，若遇空子树则创建新节点；到达叶节点时标记为叶并存 j\_value。即使管理的键数较多时，树高仍受限于键的位数，故访问步数为 \(O(\log U)\)。

\begin{algorithm}[htbp]
  \caption{Local Query Router.insert}
  \label{alg:lqr-insert}
  \begin{algorithmic}[1]
    \Function{insert}{$\mathit{key}$, $j$}
      \State $\mathit{cur}\gets\mathit{root}$
      \For{$\mathit{key}$ 的每一位 $\mathit{bit}$}
        \If{$\mathit{bit}=0$}
          \If{$\mathit{cur.child0}$ 为空} 创建子节点 \EndIf
          \State $\mathit{cur}\gets\mathit{cur.child0}$
        \Else
          \If{$\mathit{cur.child1}$ 为空} 创建子节点 \EndIf
          \State $\mathit{cur}\gets\mathit{cur.child1}$
        \EndIf
      \EndFor
      \State $\mathit{cur.is\_leaf}\gets\mathit{true}$；$\mathit{cur.j\_value}\gets j$
    \EndFunction
  \end{algorithmic}
\end{algorithm}

\begin{algorithm}[htbp]
  \caption{Local Query Router.query}
  \label{alg:lqr-query}
  \begin{algorithmic}[1]
    \Function{query}{$\mathit{key}$}
      \State $\mathit{cur}\gets\mathit{root}$
      \For{$\mathit{key}$ 的每一位 $\mathit{bit}$}
        \If{$\mathit{bit}=0$}
          \If{$\mathit{cur.child0}$ 为空} 返回 NOT\_FOUND \EndIf
          \State $\mathit{cur}\gets\mathit{cur.child0}$
        \Else
          \If{$\mathit{cur.child1}$ 为空} 返回 NOT\_FOUND \EndIf
          \State $\mathit{cur}\gets\mathit{cur.child1}$
        \EndIf
      \EndFor
      \If{$\mathit{cur.is\_leaf}$} 返回 $\mathit{cur.j\_value}$
      \Else  返回 NOT\_FOUND
      \EndIf
    \EndFunction
  \end{algorithmic}
\end{algorithm}

\begin{algorithm}[htbp]
  \caption{Local Query Router.delete}
  \label{alg:lqr-delete}
  \begin{algorithmic}[1]
    \Function{delete}{$\mathit{key}$}
      \State 沿 $\mathit{key}$ 的比特路径找到叶节点
      \State 将叶标记为非叶（$\mathit{is\_leaf}\gets\mathit{false}$），清除 $\mathit{j\_value}$
    \EndFunction
  \end{algorithmic}
\end{algorithm}

\subsection{复杂度与 Patricia 压缩}

对键长 \(\log U=\Theta(\log n)\) 的比特串，query 与 insert 沿树向下至多 \(\log U\) 步；Patricia 风格的前缀压缩\cite{morrison1968,pagh2008retrieval,lehmann2025combined} 将仅有一位不同的键合并为同一压缩边，使实测 router 深度（router\_steps\_max）远小于键长上界。每条映射存 \(j\le\log K=O(\log\log n)\) 比特，与整体 \(O(\log^{(k)}n)\) 空间目标一致\cite{bender2022otsh,kuszmaul2024filters}。

动态维护。 k-kick 踢出链中，被搬移键须 delete 旧映射再 insert 新 \(j\)；删除操作清空叶 j\_value 但不破坏共享前缀路径上的内部节点，均摊开销与 \(A[i]\) 所管理的局部键数相关。在 \(K=\mathrm{polylog}\,n\) 且 Cubby 容量 \(|I|\le K\) 的设定下，单个 \(A[i]\) 的规模受控，router 更新与查询均摊为常数时间\cite{bender2022otsh,bender2024sosa}。

\section{商压缩与键恢复}

\subsection{设计动机与信息分工}

当前层次结构中，查询路由需用映射后二进制低位定位 Facility 与 Cubby 内数组；若槽位仍存完整键，则寻址路径上的信息上下文被重复占用。商压缩借助随机可逆置换 \(\pi\)，将键映射为 \(\pi(x)\)，并按下述原则拆分信息：\(\pi(x)\) 的低 \(\log N\) 位 \(g^\ast(x)=\pi(x)[1,\log N]\) 的各段位由全局路由与局部地址上下文（Facility 编号、Cubby 容量、实际槽位下标）及 MiniArray \(B\) 中的差分元数据共同承担；槽位 storage 仅保存无法由上述上下文直接推出的高位商，即 payload
\[
  \mathrm{payload}=\pi(x)[\log N+1,\log U].
\]
查询时 Local Query Router 给出探测位置序列 j , \(h_j(x)\) 定位到 storage 中的具体位置，然后从 storage\([h_j(x)]\)、\(B[h_j(x)]\) 与 Facility/Cubby 上下文联立恢复完整 \(\pi(x)\)，再经 \(\pi^{-1}\) 得到原始键。该分工使每个元素在槽位层面仅占用 \(\log U-\log N\) 比特的有效载荷，而将路由相关的 \(\log N\) 位分散到地址结构与变长元数据中，与整体 \(O(\log^{(k)}n)\) 空间目标一致\cite{bender2022otsh,pagh2008retrieval}。

\subsection{\texorpdfstring{\(\pi(x)\) 与 \(g^\ast(x)\) 的位段划分}{pi(x) 与 g*(x) 的位段划分}}

记 \(x[a,b]\) 为 \(x\) 自第 \(a\) 位至第 \(b\) 位（自低位向高位）组成的子串。对键 \(x\) 经置换得 \(\pi(x)\)，全局路由哈希为
\[
  g^\ast(x)=\pi(x)[1,\log N].
\]
在 \(g^\ast(x)\) 内部，自高向低继续划分为三段（见图~\ref{fig:gx-layout}，第二章已给出示意）：
\begin{itemize}
  \item Facility 编号段：\(r=g^\ast(x)[\log K+1,\log N]=\lfloor g^\ast(x)/K\rfloor\)，用于定位键所属 Facility；
  \item Cubby 中间段：\(g^\ast(x)[\log|I|+1,\log K]\)，用于在 Facility 内区分不同 Cubby 及局部桶范围；
  \item Local Query Router 段：\(g_I(x)=g^\ast(x)[1,\log|I|]\in[0,|I|-1]\)，与 k-kick 偏好序列无关，仅决定键归属 \(A[i]\) 的桶号 \(i\)。
\end{itemize}

k-kick 插入可能将元素 \(x\) 写入槽位 \(i\neq g_I(x)\) 的实际位置。此时 \(g_I(x)\) 无法由槽位下标 \(i\) 直接读出，须额外保存差分
\[
  \delta=g_I(x)-i,
\]
存入 MiniArray \(B[i]\)。

\subsection{MiniArray \texorpdfstring{\(B\)}{B} 中的 MetaEntry}

每个占用槽位 \(i\) 对应一条变长元数据 MetaEntry，经 MiniArray \(B\) 的 Access/Update 接口读写。逻辑字段如下：
\begin{itemize}
  \item delta：有符号整数 \(\delta=g_I(x)-i\)，编码为 zigzag varint，适应 k-kick 搬移后槽位与 router 桶号的不一致；
  \item mid\_bits：\(g^\ast(x)\) 在 Cubby 范围内的中间位，即 \(g_k=g^\ast(x)\bmod K\) 满足 \(g_k=mid\_bits\cdot|I|+g_I(x)\) 的商部分 \(\lfloor g_k/|I|\rfloor\)；位宽 \(\texttt{mid\_w}=\lceil\log_2((K-1)/|I|+1)\rceil\)，由 \((K,|I|)\) 派生；
\end{itemize}

编码顺序为：先写入 delta 的 varint，再写入 mid\_bits 的固定位宽字段；解码时按相同布局与 MetaCodecLayout 逆序读出。\(B[i]\) 的 bitlen 随条目实际长度变化，由 Facility 级 MiniArray 的 Rank/Select 机制定位，与第二章所述变长元数据维护方式一致。

\subsection{恢复：重建 \texorpdfstring{\(\pi(x)\)}{pi(x)}}

给定 Facility 编号 \(r\)、Cubby 容量 \(|I|\)、槽位下标 \(i\)、storage\([i]\) 中的 payload 及 \(B[i]\) 解码得到的 MetaEntry，按下式自底向上重建 \(g^\ast(x)\) 与 \(\pi(x)\)：
\begin{align}
  g_I(x) &= i+\delta, \\
  g_k &= \texttt{mid\_bits}\cdot|I|+g_I(x), \\
  g^\ast(x) &= r\cdot K+g_k, \\
  \pi(x) &= (\mathrm{payload}\ll\log N)\;\vert\; g^\ast(x).
\end{align}
各步均须校验 \(g_I(x)\in[0,|I|)\)、\(g_k<K\) 及 \(g^\ast(x)<N\)；任一不满足则恢复失败，视为未命中。原始键由 \(x=\pi^{-1}(\pi(x))\) 得到。

\begin{algorithm}[htbp]
  \caption{ReconstructGxPi：从槽位上下文恢复 \(\pi(x)\)}
  \label{alg:reconstruct-gx}
  \begin{algorithmic}[1]
    \Function{ReconstructGxPi}{$r$, $N$, $K$, $|I|$, slot, payload, meta, logN}
      \State $\mathit{g\_I}\gets \mathit{slot}+\mathit{meta.delta}$
      \If{$\mathit{g\_I}\notin[0,|I|)$} \State 返回失败 \EndIf
      \State $\mathit{g\_k}\gets \mathit{meta.mid\_bits}\cdot|I|+\mathit{g\_I}$
      \If{$\mathit{g\_k}\ge K$} \State 返回失败 \EndIf
      \State $\mathit{g\_star}\gets r\cdot K+\mathit{g\_k}$
      \If{$\mathit{g\_star}\ge N$} \State 返回失败 \EndIf
      \State 返回 $(\mathit{payload}\ll\mathit{logN})\;\vert\;\mathit{g\_star}$
    \EndFunction
  \end{algorithmic}
\end{algorithm}

\subsection{查询校验}

查询键 \(x\) 时，Router 给出探测下标 \(j\)，访问候选槽位 \(i=h_j(x)\)。只有当完整恢复的值与查询的值相同，才认为查询命中：

\begin{algorithm}[htbp]
  \caption{SlotPiMatches：槽位命中校验}
  \label{alg:slot-match}
  \begin{algorithmic}[1]
    \Function{SlotPiMatches}{cubby, slot, $r$, table, query\_gx}
      \If{cubby.slots[slot] 为空} \State 返回 false \EndIf
      \State meta $\gets$ \Call{DecodeMetaEntry}{cubby.B[slot]}
      \State gx $\gets$ \Call{ReconstructGxPi}{$r$, $N$, $K$, $|I|$, slot, payload, meta, logN}
      \If{gx 失败} \State 返回 false \EndIf
      \State 返回 gx $=$ query\_gx
    \EndFunction
  \end{algorithmic}
\end{algorithm}

Router 仅负责将查询映射到 \(O(I)\) 个候选位置；最终是否命中必须由 payload、\(B\) 与 \((r,N,K,|I|)\) 联立恢复的 \(\pi(x)\) 与查询侧计算的 \(\pi(x)\) 一致方可确认。该“Router 定位 + 商压缩校验”的两段式路径，是系统能够不存储完整键的前提\cite{bender2022otsh,pagh2008retrieval}。


\subsection{空间性质}

单个占用槽位的比特开销可分解为：（1）storage 中 \(\log U-\log N\) 比特 payload；（2）\(B[i]\) 中变长 MetaEntry，其中， \(\delta\) 采用 varint 编码，长度会随 \(|g_I(x)-i|\) 增大而增加，mid\_bits 占 \(\Theta(\log(K/|I|))\) 比特，router\_bits 与 insert\_bits 为常数宽度。Facility 编号 \(r\) 与 Cubby 容量 \(|I|\) 作为查询路径上的上下文信息，不重复写入每个槽位。整体上来看，商压缩将 \(\pi(x)\) 中与寻址重叠的低位从槽位载荷中抽离，使附加空间与 \(O\log^{(k)}n\) 级预算相容；MiniArray \(B\) 的变长编码则使元数据按条目实际所需位数记录，避免为所有槽位预留固定最大键长\cite{kurpicz2023pachash,lehmann2024shockhash}。

\section{插入、查询与删除}

插入。 计算 \(g^\ast(x)=\pi(x)[1,\log N]\) 与 Facility 编号 \(r\)；在对应 Facility 的 tail Cubby 内执行 InsertKick（算法~\ref{alg:insert-kick}），写入 storage、更新 \(M\) 与 \(B[i]\)，并在 \(A[g_I(x)]\) 中 insert(\(\pi(x),j\))。tail 满则挂入 1-tiered cubby 并新建 tail，必要时按算法~\ref{alg:tier-schedule} 触发合并\cite{bender2022otsh,mitzenmacher2001}。

查询。 计算 \(g^\ast(x)\) 定位 Facility 与 Cubby；令 \(i=g_I(x)\)，在 \(A[i]\) 中 query(\(\pi(x)\)) 得 \(j\)；访问 \(h_j(x)\)，用 payload、\(B\)、\(r\) 重建 \(g^\ast(x)\) 并与 \(\pi(x)\) 比较。该路径仅含常数次 MiniArray 与 Local Query Router 访问，查询时间为 \(O(1)\)\cite{bender2022otsh}。

删除。 定位元素所在 Cubby 与槽位，清空 storage，更新 \(A[g_I(x)]\).delete(\(\pi(x)\))、\(M\)与 \(B\)。若删除发生在非 tail Cubby，按算法~\ref{alg:tail-ops} 从 tail 随机抽取元素回填；回填写入同样走 k-kick 与 Local Query Router 更新路径，以恢复各 bin 的饱和与深度不变量。

表级扩缩容时，\(N\)、\(K\)、\(r_j\) 与商压缩上下文变化，须先恢复原始键再按新参数重插，而不能直接复制旧槽位的 payload 与 \(B\)\cite{bender2022otsh,bender2024sosa}。

\section{本章小结}

本章给出分层哈希表的主要结构与算法：Facility 级 MiniArray 以 B 树、Rank/Select 与 Split/Merge 维护变长元数据\cite{bayer1972,jacobson1989}；多层级 Cubby 与 tail 机制定义 tier 调度与删除回填不变量\cite{bender2024iceberg,pandey2023iceberght}；k-kick 经偏好序列、饱和判定、InsertKick 与 ReInsert 实现 \(O(k)\) 交换插入\cite{bender2022otsh,pagh2004cuckoo}；Local Query Router 经 \(A[i]\) 前缀树将查询降为 \(O(1)\)\cite{morrison1968,pagh2008retrieval}。四组件同步更新规则贯穿插入、踢出、删除回填与扩缩容重插；第四章报告各参数下的性能测量。

\chapter{系统性能实验}
\label{ch:exp}

\section{实验目的与环境}

本章基于第二、三章给出的结构与参数设定，对已实现的哈希系统进行性能测试和功能验证。实验围绕操作延迟、k-kick 踢出次数、多层级 Cubby 向下拆分（rebuild\_down）与向上合并（rebuild\_up）三类指标展开，用于检查第三章各模块在实际运行中的效果\cite{bender2022otsh,bender2024sosa}。

实验的系统代码基于 C++23 标准开发实现，在 AMD Ryzen 7 6800H 处理器、Windows 10 环境下运行。元素规模 \(n\in\{10^3,10^4,10^5\}\)，Facility 规模 \(K\) 按分组取多档（基准 \(K=\log^3 n\)），k-kick 最大深度 \(k\in\{3,4,5\}\)（基准 \(k=4\)），NODE\_MAX\_BITS 常数 \(c=2\)。键由 std::mt19937\_64 在固定种子 seed=1（经 seed1/seed2/seed3 派生表内哈希参数）下均匀生成；实验报告单次 workload 内各操作类型的平均延迟。

\section{实验方案}

\subsection{实验分组与工作量}

实验分为四组，如表~\ref{tab:exp-plan} 所示。规模 \(n\) 组与 \(K\) 组、\(k\) 组采用键均匀分布于各 Facility 的混合增删负载；Cubby 动态重构组按纯插入、纯删除及不同插入/删除比例构造 workload，分别触发 rebuild\_up 与 rebuild\_down。

\begin{table}[htbp]
  \centering
  \caption{实验分组与主要验证内容}
  \label{tab:exp-plan}
  \small
  \begin{tabular}{p{1.6cm}p{3.6cm}p{5.8cm}}
    \toprule
    分组 & 负载方式 & 主要验证内容 \\
    \midrule
    规模 \(n\) & 均匀分布、混合增删 & 操作延迟、踢出、rebuild \\
    Facility \(K\) & 均匀分布、混合增删 & 操作延迟、踢出 \\
    k-kick & 均匀分布、高负载 & 操作延迟、踢出 \\
    Cubby & Insert/Delete/混合比例 & rebuild\_down、rebuild\_up、重构耗时 \\
    \bottomrule
  \end{tabular}
\end{table}

\subsection{评价指标}

为量化评估系统的基础性能，实验记录并分析以下指标：

\begin{itemize}
  \item 操作延迟：插入、命中查询和删除的平均延迟。该指标用于衡量哈希 系统基础操作的运行性能；
  \item k-kick：workload 内累计踢出次数，以及单次插入/删除搬移最大次数 insert\_moved\_max、delete\_moved\_max。该指标用于观察 k-kick 核心模块的具体表现，并判断模块的基础踢出功能是否生效；
  \item j-tiered Cubby：rebuild\_down（向下拆分）与 rebuild\_up（向上合并）触发次数。该指标用于验证多层级的 Cubby 的合并与拆分功能是否正常触发；
\end{itemize}

\section{规模 \texorpdfstring{$n$}{n} 实验}

固定 \(K=\log^3 n\)、\(k=4\)、\(c=2\)，令 \(n\) 从 \(10^3\) 增至 \(10^5\)。表~\ref{tab:exp-n-fixed} 给出各档测量结果。

\begin{table}[htbp]
  \centering
  \caption{规模 \(n\) 实验结果（\(K=\log^3 n\)，\(k=4\)，键均匀分布）}
  \label{tab:exp-n-fixed}
  \footnotesize
  \begin{tabular}{lrrrrrr}
    \toprule
    $n$ &
    ins ($\mu$s) & qry ($\mu$s) & del ($\mu$s) &
    总踢出 & rebuild\_down & rebuild\_up \\
    \midrule
    $10^3$  & 49.1 & 2.2  & 66.0 & 80     & 11    & 179   \\
    $10^4$  & 42.0 & 2.7  & 71.6 & 1157   & 77    & 930   \\
    $10^5$  & 60.8 & 5.1  & 2051.9 & 10485 & 11447 & 21502 \\
    \bottomrule
  \end{tabular}
\end{table}

\noindent 表中 rebuild\_down 对应向下拆分的次数，rebuild\_up 对应向上合并的次数。

\begin{figure}[htbp]
  \centering
  \begin{tikzpicture}
    \begin{axis}[
      width=0.88\linewidth,
      height=6.2cm,
      xmode=log,
      log basis x={10},
      xtick={1000,10000,100000},
      xticklabels={$10^3$,$10^4$,$10^5$},
      xmin=800, xmax=120000,
      xlabel={规模 $n$},
      ylabel={平均延迟 ($\mu$s)},
      legend style={font=\small},
      legend pos=north west,
      grid=major,
      grid style={dashed,gray!30},
    ]
    \addplot[mark=*, blue, thick] coordinates {(1000,49.1) (10000,42.0) (100000,60.8)};
    \addlegendentry{插入}
    \addplot[mark=square*, red, thick] coordinates {(1000,2.2) (10000,2.7) (100000,5.1)};
    \addlegendentry{查询}
    \addplot[mark=triangle*, green!60!black, thick] coordinates {(1000,66.0) (10000,71.6) (100000,2051.9)};
    \addlegendentry{删除}
    \end{axis}
  \end{tikzpicture}
  \caption{规模 $n$ 实验中三种操作平均延迟（$K=\log^3 n$，$k=4$，键均匀分布）}
  \label{fig:exp-n-latency}
\end{figure}

随 \(n\) 增大，命中查询延迟由 2.2\,$\mu$s 升至 5.1\,$\mu$s，相对三个数量级跨度增幅平缓，与 Local Query Router 的常数步定位一致\cite{bender2022otsh,morrison1968}。插入延迟在 42.0--60.8\,$\mu$s 之间，\(n=10^3\) 略高（49.1\,$\mu$s），与 Facility 数较少、单桶局部竞争有关。删除延迟由 66.0\,$\mu$s 增至 2051.9\,$\mu$s，增幅远高于插入：\(n=10^5\) 时 rebuild 合计约 3.3 万次，删除路径上的 tail 回填、tier 拆分及 MiniArray 同步均叠加在热路径上（见图~\ref{fig:exp-n-latency}）。累计踢出由 80 增至 \(1.0\times 10^4\)，rebuild\_down 由 11 增至 11447 次、rebuild\_up 由 179 增至 21502 次；大规模下 j-tiered Cubby 维护最为频繁\cite{bender2024iceberg,pandey2023iceberght}。

\section{Facility 规模 \texorpdfstring{$K$}{K} 实验}

固定 \(k=4\)，在 \(n\in\{10^3,10^4,10^5\}\) 三档下各取多组 \(K\) 比较三操作延迟与踢出。表~\ref{tab:exp-K-1e3}--\ref{tab:exp-K-1e5} 汇总实测结果（键均匀分布）。

\begin{table}[htbp]
  \centering
  \caption{Facility 规模 \(K\) 实验（\(n=10^3\)，\(k=4\)）}
  \label{tab:exp-K-1e3}
  \small
  \begin{tabular}{lrrrr}
    \toprule
    $K$ & ins ($\mu$s) & qry ($\mu$s) & del ($\mu$s) & 总踢出 \\
    \midrule
    16   & 22.6 & 2.1 & 52.8 & 112 \\
    64   & 49.9 & 2.0 & 67.7 & 60  \\
    256  & 51.9 & 1.9 & 66.0 & 58  \\
    \bottomrule
  \end{tabular}
\end{table}

\begin{table}[htbp]
  \centering
  \caption{Facility 规模 \(K\) 实验（\(n=10^4\)，\(k=4\)）}
  \label{tab:exp-K-fixed}
  \small
  \begin{tabular}{lrrrr}
    \toprule
    $K$ & ins ($\mu$s) & qry ($\mu$s) & del ($\mu$s) & 总踢出 \\
    \midrule
    64    & 56.9 & 2.8 & 75.9 & 813 \\
    256   & 54.1 & 2.7 & 70.9 & 916 \\
    1024  & 52.9 & 2.6 & 72.0 & 963 \\
    \bottomrule
  \end{tabular}
\end{table}

\begin{table}[htbp]
  \centering
  \caption{Facility 规模 \(K\) 实验（\(n=10^5\)，\(k=4\)）}
  \label{tab:exp-K-1e5}
  \small
  \begin{tabular}{lrrrr}
    \toprule
    $K$ & ins ($\mu$s) & qry ($\mu$s) & del ($\mu$s) & 总踢出 \\
    \midrule
    256   & 98.9  & 4.2 & 504.9  & 6889 \\
    1024  & 185.2 & 3.4 & 518.3  & 6883 \\
    4096  & 64.9  & 3.1 & 127.1  & 8756 \\
    \bottomrule
  \end{tabular}
\end{table}

\(n=10^3\) 时 \(K=16\) 插入最快（22.6\,$\mu$s），但踢出最多（112 次）；\(K\ge 64\) 时踢出降至约 60 次，增删延迟略升，体现小 \(n\) 下 Facility 过细反而增加路由与维护开销。\(n=10^4\) 三档 \(K\) 下查询均在 2.6--2.8\,$\mu$s，踢出随 \(K\) 增大变化不大（813--963），说明中等规模下 \(K\) 对均匀负载的边际影响有限。\(n=10^5\) 时 \(K=4096\) 删除延迟（127.1\,$\mu$s）明显低于 \(K=256/1024\)（504--518\,$\mu$s），与大规模 rebuild 主要落在较小 \(K\) 档、大 Facility 缓解局部冲突一致；\(K=1024\) 插入异常偏高（185.2\,$\mu$s）可能与 Facility 数与 Cubby 布局的组合有关，部署时宜结合目标 \(n\) 与负载率选取 \(K\)\cite{bender2022otsh,bender2024sosa}。

\section{k-kick 深度 \texorpdfstring{$k$}{k} 实验}

对 \(k\in\{3,4,5\}\) 在各规模下执行相同高负载混合增删 workload：\(n=10^3\) 取 \(K=64\)，\(n=10^4\) 取 \(K=256\)，\(n=10^5\) 取 \(K=4096\)。表~\ref{tab:exp-k-1e3}--\ref{tab:exp-k-1e5} 给出完整结果；图~\ref{fig:exp-k-bpk} 汇总三档规模下 bits/key 随 \(k\) 的变化。

\begin{table}[htbp]
  \centering
  \caption{k-kick 深度 \(k\) 实验（\(n=10^3\)，\(K=64\)，高负载）}
  \label{tab:exp-k-1e3}
  \small
  \begin{tabular}{lrrrr}
    \toprule
    $k$ & ins ($\mu$s) & qry ($\mu$s) & del ($\mu$s) & 总踢出 \\
    \midrule
    3 & 71.2 & 2.1 & 85.8 & 61  \\
    4 & 73.0 & 2.1 & 88.7 & 104 \\
    5 & 77.5 & 2.3 & 97.7 & 194 \\
    \bottomrule
  \end{tabular}
\end{table}

\begin{table}[htbp]
  \centering
  \caption{k-kick 深度 \(k\) 实验（\(n=10^4\)，\(K=256\)，高负载）}
  \label{tab:exp-k-fixed}
  \small
  \begin{tabular}{lrrrr}
    \toprule
    $k$ & ins ($\mu$s) & qry ($\mu$s) & del ($\mu$s) & 总踢出 \\
    \midrule
    3 & 54.5 & 2.7 & 72.6 & 265  \\
    4 & 55.3 & 2.6 & 72.6 & 900  \\
    5 & 56.4 & 2.5 & 70.2 & 1587 \\
    \bottomrule
  \end{tabular}
\end{table}

\begin{table}[htbp]
  \centering
  \caption{k-kick 深度 \(k\) 实验（\(n=10^5\)，\(K=4096\)，高负载）}
  \label{tab:exp-k-1e5}
  \small
  \begin{tabular}{lrrrr}
    \toprule
    $k$ & ins ($\mu$s) & qry ($\mu$s) & del ($\mu$s) & 总踢出 \\
    \midrule
    3 & 60.3 & 3.1 & 2229.4 & 2772  \\
    4 & 62.6 & 3.2 & 2003.9 & 8707  \\
    5 & 69.1 & 3.3 & 1734.3 & 14986 \\
    \bottomrule
  \end{tabular}
\end{table}


\begin{figure}[htbp]
  \centering
  \begin{tikzpicture}
    \begin{axis}[
      width=0.88\linewidth,
      height=6.2cm,
      xmin=2.5, xmax=5.5,
      xtick={3,4,5},
      xlabel={k-kick 深度 $k$},
      ylabel={bits/key},
      legend style={font=\small},
      legend pos=north east,
      grid=major,
      grid style={dashed,gray!30},
      ymin=184, ymax=200,
    ]
    \addplot[mark=*, blue, thick] coordinates {(3,191.19) (4,189.67) (5,186.53)};
    \addlegendentry{$n=10^3$}
    \addplot[mark=square*, red, thick] coordinates {(3,192.20) (4,189.10) (5,186.12)};
    \addlegendentry{$n=10^4$}
    \addplot[mark=triangle*, green!60!black, thick] coordinates {(3,198.39) (4,195.48) (5,192.65)};
    \addlegendentry{$n=10^5$}
    \end{axis}
  \end{tikzpicture}
  \caption{k-kick 深度 $k$ 实验中 bits/key}
  \label{fig:exp-k-bpk}
\end{figure}

\noindent insert\_moved\_max 与 delete\_moved\_max 各档均不超过 \(k\)。

三档规模下累计踢出均随 \(k\) 增大而上升：\(n=10^3\) 由 61 增至 194，\(n=10^4\) 由 265 增至 1587，\(n=10^5\) 由 2772 增至 14986\cite{bender2022otsh,pagh2004cuckoo}。查询延迟在各规模均不超过 3.3\,$\mu$s，kick 深度主要作用于增删路径。\(n=10^5\) 删除延迟达毫秒级，且 \(k=5\) 时（1734.3\,$\mu$s）低于 \(k=3/4\)，与大规模 rebuild 叠加、较深 kick 链在特定负载下分摊搬移的实测表现一致；中小规模下增删延迟随 \(k\) 变化相对平缓\cite{bender2022otsh,bender2024sosa}。图~\ref{fig:exp-k-bpk},表明，随着 k-kick 树深度的增加，每键占用的内存呈现下降趋势，这与论文核心观点相符。

\section{多层级 Cubby 动态重构实验}

固定 \(k=4\)，在 \(n\in\{10^3,10^4,10^5\}\) 下分别取 \(K=64\)、\(256\)、\(4096\)，按表~\ref{tab:exp-tier-1e3}--\ref{tab:exp-tier-1e5} 四种 workload 执行操作序列，统计 rebuild 次数与单次重构平均耗时。

\begin{table}[htbp]
  \centering
  \caption{多层级 Cubby 动态重构（\(n=10^3\)，\(K=64\)）}
  \label{tab:exp-tier-1e3}
  \small
  \begin{tabular}{lrrr}
    \toprule
    Workload      & rebuild\_down & rebuild\_up & 平均重构耗时 (ms) \\
    \midrule
    Insert-only   & 1   & 31  & 3.72 \\
    Delete-only   & 28  & 0   & 6.24 \\
    80/20 混合    & 0   & 24  & 3.71 \\
    50/50 混合    & 6   & 25  & 3.82 \\
    \bottomrule
  \end{tabular}
\end{table}

\begin{table}[htbp]
  \centering
  \caption{多层级 Cubby 动态重构（\(n=10^4\)，\(K=256\)）}
  \label{tab:exp-tier-fixed}
  \small
  \begin{tabular}{lrrr}
    \toprule
    Workload      & rebuild\_down & rebuild\_up & 平均重构耗时 (ms) \\
    \midrule
    Insert-only   & 0   & 384 & 3.47 \\
    Delete-only   & 384 & 2   & 6.10 \\
    80/20 混合    & 8   & 355 & 3.32 \\
    50/50 混合    & 68  & 377 & 3.35 \\
    \bottomrule
  \end{tabular}
\end{table}

\begin{table}[htbp]
  \centering
  \caption{多层级 Cubby 动态重构（\(n=10^5\)，\(K=4096\)）}
  \label{tab:exp-tier-1e5}
  \small
  \begin{tabular}{lrrr}
    \toprule
    Workload      & rebuild\_down & rebuild\_up & 平均重构耗时 (ms) \\
    \midrule
    Insert-only   & 0    & 3591 & 5.64 \\
    Delete-only   & 4211 & 621  & 9.76 \\
    80/20 混合    & 177  & 3588 & 5.52 \\
    50/50 混合    & 375  & 3504 & 5.13 \\
    \bottomrule
  \end{tabular}
\end{table}

三档规模下规律一致：纯插入以 rebuild\_up 为主（\(n=10^3\) 为 31 次，\(n=10^5\) 为 3591 次），纯删除以 rebuild\_down 为主（28--4211 次），对应 Cubby 随元素增减而向上合并或向下拆分\cite{bender2024iceberg,pandey2023iceberght}。80/20 混合负载在各规模均以 rebuild\_up 为主；50/50 混合时两方向均显著（如 \(n=10^4\) 为 68/377，\(n=10^5\) 为 375/3504），删除比例升高使 rebuild\_down 占比上升。平均重构耗时由 3.3--6.2\,ms（\(n\le 10^4\)）增至 5.1--9.8\,ms（\(n=10^5\)），仍远低于大规模混合 workload 中单次删除的均摊延迟。

\section{本章小结}

本章在 \(n\)、Facility 规模 \(K\)、k-kick 深度 \(k\) 及 Cubby 动态重构四组设定下，于 \(n\in\{10^3,10^4,10^5\}\) 完成测量。表~\ref{tab:exp-n-fixed} 与图~\ref{fig:exp-n-latency} 显示查询延迟由 2.2\,$\mu$s 升至 5.1\,$\mu$s，\(n=10^5\) 删除延迟达 2051.9\,$\mu$s；\(K\)、\(k\) 与 Cubby 三组实验在表~\ref{tab:exp-K-1e3}--\ref{tab:exp-tier-1e5} 中给出三档规模对照，insert\_moved\_max 各档均不超过 \(k\)，纯插入/纯删除 workload 分别主导 rebuild\_up 与 rebuild\_down。上述结果与 Local Query Router 常数步查询设计及 k-kick 交换上界在实测中可对照。
\chapter{冷热分层存储对照实验}
\label{ch:hotcold}

\section{实验目的与方案设置}

第四章验证的是基础多层级结构中 facility、k-kick、局部路由和重构模块能否正常工作。第三章已经说明冷热分层键存储的实现方式：插入的键先进入内存 tail cubby，tail cubby 容量满，切换时再批量下沉到 SQLite3 数据库存储。本章在该实现基础上设置三组对照实验，观察内存 tail cubby 与 SQLite3 冷数据层结合后带来的系统性能效果。第五章与第四章延迟数值不可直接对比，是因为测试的数据方式不同，属于合理范围。

实验设置了三个对照方案：(1) 原始的多层级 cubby 结构，数据均保存在内存多层级 cubby 中；(2) 纯数据库结构，所有插入、查询和删除都直接访问 SQLite3；(3) 冷热分层结构，近期写入的数据先进入内存 tail cubby，tail cubby 写满后再批量写入 SQLite3。在相同的数据规模和相同的测试数据情况下进行实验，避免数据层面的差异带来的误差。

\section{负载与统计口径}

实验选取 \(5\times10^3\)、\(10^4\)、\(5\times10^4\)、\(10^5\)、\(5\times10^5\) 五档规模。每一档先插入全部键，再进行同规模次数的查询，随后删除近期的 20\% 键。测试数据采用 C++ 标准库的 std::mt19937\_64 随机数生成器产生互不重复的 64 位整数键，以保证键值分布的均匀性和实验结果的可复现性。完全插入系统后，查询阶段按照键的插入顺序建立访问排名，其中最新插入的键具有更高访问概率，并依据 Zipf 分布进行采样，以模拟具有时间局部性的访问负载\cite{cooper2010ycsb}。删除阶段选择最近插入的 20\% 键，用于模拟查询与删除主要针对近期写入的键。

统计指标包括插入平均延迟、查询平均延迟、删除平均延迟、平均每键空间占用和正确率。平均延迟只统计对应操作本身，不把收尾阶段的数据库刷盘和检查过程计入单次操作。平均每键空间占用按删除后的剩余键数计算，内存结构统计逻辑元数据位数，数据库结构统计 SQLite3 主数据库文件大小。

\section{多层级内存结构对照}

表~\ref{tab:hotcold-memory-latency} 对比了多层级内存结构与冷热分层结构的实验结果。多层级内存结构的查询延迟在五档规模下均低于 1\,$\mu$s，说明第四章中的局部路由在新增负载下仍然保持较短查询路径。冷热分层结构的查询延迟为 0.399--0.550\,$\mu$s，除小规模 \(5\times10^3\) 外，与多层级内存结构处于同一数量级。

\begin{table}[htbp]
  \centering
  \caption{多层级内存结构与冷热分层结构延迟对照}
  \label{tab:hotcold-memory-latency}
  \small
  \begin{tabular}{llrrr}
    \toprule
    规模 & 方案 & 插入 ($\mu$s) & 查询 ($\mu$s) & 删除 ($\mu$s) \\
    \midrule
    \(5\times10^3\)   & 多层级内存结构 & 42.189 & 0.335 & 140.200 \\
    \(5\times10^3\)   & 冷热分层结构   & 10.231 & 0.476 & 8.605 \\
    \(10^4\)          & 多层级内存结构 & 41.804 & 0.345 & 123.759 \\
    \(10^4\)          & 冷热分层结构   & 9.686  & 0.399 & 11.341 \\
    \(5\times10^4\)   & 多层级内存结构 & 57.381 & 0.510 & 181.186 \\
    \(5\times10^4\)   & 冷热分层结构   & 17.198 & 0.464 & 11.148 \\
    \(10^5\)          & 多层级内存结构 & 61.702 & 0.516 & 167.669 \\
    \(10^5\)          & 冷热分层结构   & 17.432 & 0.486 & 12.638 \\
    \(5\times10^5\)   & 多层级内存结构 & 83.232 & 0.708 & 129.434 \\
    \(5\times10^5\)   & 冷热分层结构   & 19.271 & 0.550 & 23.660 \\
    \bottomrule
  \end{tabular}
\end{table}

插入和删除结果更能体现两类结构的差异。多层级内存结构需要维护 tail cubby 回填以及层级合并、拆分，删除延迟在五档规模下为 123.759--181.186\,$\mu$s，\(5\times10^3\) 时为 140.200\,$\mu$s。冷热分层结构删除近期键时多发生在内存 tail cubby 或尚未完全下沉的数据区域，删除延迟为 8.605--23.660\,$\mu$s。在 Zipf 写后读负载下，冷热分层插入和删除延迟更低，删除延迟比多层级内存结构低一个数量级。

\section{纯数据库结构对照}

表~\ref{tab:hotcold-sqlite-latency} 对比了纯数据库结构与冷热分层结构的实验结果。纯数据库结构的查询都需要经过 SQLite3\cite{gaffney2022sqlite}，五档规模下查询延迟为 1.190--1.433\,$\mu$s；冷热分层结构在相同查询序列下为 0.399--0.550\,$\mu$s。该差异来自负载中的近期访问特征：一部分新近写入的键仍在内存 tail cubby 中，查询不必进入数据库。

\begin{table}[htbp]
  \centering
  \caption{纯数据库结构与冷热分层结构延迟对照}
  \label{tab:hotcold-sqlite-latency}
  \small
  \begin{tabular}{llrrr}
    \toprule
    规模 & 方案 & 插入 ($\mu$s) & 查询 ($\mu$s) & 删除 ($\mu$s) \\
    \midrule
    \(5\times10^3\)   & 纯数据库结构 & 12.064 & 1.190 & 13.575 \\
    \(5\times10^3\)   & 冷热分层结构 & 10.231 & 0.476 & 8.605 \\
    \(10^4\)          & 纯数据库结构 & 12.639 & 1.257 & 13.735 \\
    \(10^4\)          & 冷热分层结构 & 9.686  & 0.399 & 11.341 \\
    \(5\times10^4\)   & 纯数据库结构 & 13.659 & 1.279 & 12.301 \\
    \(5\times10^4\)   & 冷热分层结构 & 17.198 & 0.464 & 11.148 \\
    \(10^5\)          & 纯数据库结构 & 13.684 & 1.284 & 13.464 \\
    \(10^5\)          & 冷热分层结构 & 17.432 & 0.486 & 12.638 \\
    \(5\times10^5\)   & 纯数据库结构 & 20.567 & 1.433 & 25.888 \\
    \(5\times10^5\)   & 冷热分层结构 & 19.271 & 0.550 & 23.660 \\
    \bottomrule
  \end{tabular}
\end{table}

插入延迟并不是冷热分层结构在所有规模下都更低，\(5\times10^4\) 与 \(10^5\) 两档中，冷热分层结构插入分别为 17.198\,$\mu$s 和 17.432\,$\mu$s，高于纯数据库结构的 13.659\,$\mu$s 和 13.684\,$\mu$s，但是这处于能接受的范围内。这与内存 tail cubby 中的 k-kick 插入以及后续批量写入有关。到 \(5\times10^5\) 时，冷热分层结构插入为 19.271\,$\mu$s，略低于纯数据库结构的 20.567\,$\mu$s。删除方面，冷热分层结构五档均低于纯数据库结构，但差距不如查询明显。因此，收益主要体现在近期查询与删除，写入并不稳定优于纯 SQLite3。

\section{空间占用与正确性对照}

表~\ref{tab:hotcold-space-correctness} 汇总三种方案的平均每键空间占用和正确率的实验结果。在本文测试用例与访问序列下，三类结构的查询结果均与预期一致，说明原型实现保持了基本插入、查询和删除语义。空间方面，多层级内存结构为 64.427--68.305 位/键，明显低于使用 SQLite3 的方案。这是因为多层级内存结构只统计逻辑元数据位数，而 SQLite3 主数据库文件包含页面、表结构和索引等额外的固定开销。

\begin{table}[htbp]
  \centering
  \caption{三种方案空间占用与正确率对照}
  \label{tab:hotcold-space-correctness}
  \small
  \begin{tabular}{llrr}
    \toprule
    规模 & 方案 & 每键位数 & 正确率（\%） \\
    \midrule
    \(5\times10^3\)   & 多层级内存结构 & 65.333  & 100.000 \\
    \(5\times10^3\)   & 纯数据库结构   & 204.800 & 100.000 \\
    \(5\times10^3\)   & 冷热分层结构   & 188.416 & 100.000 \\
    \(10^4\)          & 多层级内存结构 & 64.427  & 100.000 \\
    \(10^4\)          & 纯数据库结构   & 200.704 & 100.000 \\
    \(10^4\)          & 冷热分层结构   & 184.320 & 100.000 \\
    \(5\times10^4\)   & 多层级内存结构 & 67.141  & 100.000 \\
    \(5\times10^4\)   & 纯数据库结构   & 160.563 & 100.000 \\
    \(5\times10^4\)   & 冷热分层结构   & 154.829 & 100.000 \\
    \(10^5\)          & 多层级内存结构 & 67.904  & 100.000 \\
    \(10^5\)          & 纯数据库结构   & 160.973 & 100.000 \\
    \(10^5\)          & 冷热分层结构   & 153.600 & 100.000 \\
    \(5\times10^5\)   & 多层级内存结构 & 68.305  & 100.000 \\
    \(5\times10^5\)   & 纯数据库结构   & 155.075 & 100.000 \\
    \(5\times10^5\)   & 冷热分层结构   & 151.880 & 100.000 \\
    \bottomrule
  \end{tabular}
\end{table}

在使用 SQLite3 的两类方案之间，冷热分层结构的每键位数略低于纯数据库结构，例如 \(10^5\) 时为 153.600 位/键，而纯数据库结构为 160.973 位/键。这一差异与部分近期数据仍保留在内存 tail cubby、数据库文件中实际写入量和页面分配情况有关。该结果不能简单解释为冷热分层结构空间上总是更优，因为内存 tail cubby 的逻辑元数据和数据库页面开销属于不同来源；但在本实验口径下，冷热分层结构没有因为增加内存热层而超过纯数据库结构的总空间统计，相反因为原本的紧凑哈希设计，存储在 tail cubby 中的数据占用更少的空间。

\section{本章小结}

本章用相同输入和统计口径比较了多层级内存结构、纯数据库结构和冷热分层结构。在具有时间局部性的查询与删除负载下，冷热分层带来的性能收益来自内存热层：近期查询较少直接访问 SQLite3，近期删除也较少触发多层级回填。写入并不总是占优，\(5\times10^4\) 和 \(10^5\) 两档下的插入延迟高于纯数据库结构。若查询目标大多落在冷层，或者 SQLite3 已经有较高页面缓存命中率，热层判断和状态维护反而可能增加成本。因此，本章结论只对应访问局部性较强的负载，不能解释为冷热分层无条件更快。

\chapter{总结与展望}
\label{ch:concl}

\section{工作总结}

本文以 Bender 等人提出的最优时间/空间权衡框架\cite{bender2022otsh} 为理论基础，实现了基于 k-kick 树的分层哈希表的结构与算法设计，并据此实现简化原型、开展性能实验。实现中的主要难点在于将 mini-array、k-kick 与 local query router 串联合并为统一系统，并在 facility--cubby--slot 分层组织下保持辅助结构同步不变量。在此基础上，本文进一步加入冷热分层键存储方案，用内存 tail cubby 存储近期数据，用 SQLite3 管理下沉后的冷数据\cite{gaffney2022sqlite,dong2021rocksdb}。

设计部分采用 facility--cubby--slot 分层组织。更新路径主要依赖 mini-array 维护变长元数据，并通过 k-kick 控制插入过程中的有限元素交换。查询路径由 local query router 定位候选槽位，再结合商化压缩信息恢复和校验键。每次发生搬移时，槽位数组、mini-array \(B\)、\(M\) 与 local query router 都需要同步更新。

实验测量分为两部分。第四章表~\ref{tab:exp-n-fixed}--\ref{tab:exp-tier-fixed} 及图~\ref{fig:exp-n-latency} 检查基础多层级结构的功能与维护成本；踢出深度变化时，单次搬移次数仍不超过对应的 \(k\)，与前文交换上界一致。第五章表~\ref{tab:hotcold-memory-latency}--\ref{tab:hotcold-space-correctness} 对比多层级内存结构、纯数据库结构和冷热分层结构。冷热分层在近期热点查询和近期删除负载下延迟较低，代价是空间占用高于纯内存紧凑结构。

本文实现了 k-kick 哈希表的简化原型，记录了与严格理论模型的差异，并给出测量数据与实验参数。冷热分层部分冷数据改为 SQLite3 记录后，近期访问更快，但失去了多层级 cubby 的紧凑性。

\section{不足与局限}

本文实验还有四方面不足：

(1) 实现的系统原型尚未与 IcebergHT\cite{pandey2023iceberght}、PaCHash\cite{kurpicz2023pachash} 等同类系统做同机对比，因此只能说明自身参数变化下的行为，难以直接评价相对性能；

(2) 未验证 k-kick 的树高，区间大小的取值等因素是否会影响系统性能效果；

(3) 第四章采用的合成负载较理想，规模覆盖也有限；

(4) 第五章更关注写后读和近期删除，对随机冷查询、数据库页面缓冲充分以及更复杂持久化策略的覆盖还不够。

\section{展望}

后续工作可在现有实验基础上进一步补充对照实验和运行指标。本文目前主要比较了基础多层级结构、纯 SQLite 结构与冷热分层结构的表现，后续可在相同硬件环境、相同数据规模和相同负载设置下，引入 PaCHash 或 IcebergHT 等紧凑哈希结构作为同机 baseline，从而更准确地判断本文方案在空间开销、查询延迟和更新代价上的相对位置\cite{kurpicz2023pachash,lehmann2026mphsurvey,pibiri2021pthash}。同时，现有实验主要关注单次操作延迟，对整体运行状态的刻画仍不充分，后续可继续记录进程级内存占用、SQLite 运行时文件大小、批量写入吞吐量以及混合负载下的查询和更新吞吐表现\cite{cooper2010ycsb}。在此基础上，还可以针对插入和删除热路径继续优化，并比较不同数据规模下 facility 规模、k-kick 深度和区间大小对性能的影响，以分析冷热分层结构在不同负载条件下是否存在更合适的参数组合。


%---------------------------------------------------------------------
%	参考文献
%---------------------------------------------------------------------

% 生成参考文献页
\printbibliography

%---------------------------------------------------------------------
%	致谢
%---------------------------------------------------------------------

\begin{acknowledgement}
  行文至此，落笔为终。感谢我的导师、感谢身边的每一个同学、感谢遇到的每一个人。凡是过往，皆为序章，很喜欢看到的一句话，人生如棋，落子无悔。愿我们跃入人海，各自灿烂。
\end{acknowledgement}

%---------------------------------------------------------------------
%	学术简历
%---------------------------------------------------------------------

% 详见手册中“成果列表”一节
% \njuchapter{学术成果}
% \section*{个人简历}
% \njupaperlist[攻读博士学位期间发表的学术论文]{preskill2018}

%---------------------------------------------------------------------
%	附录部分
%---------------------------------------------------------------------

% 附录部分使用单独的字母序号
\appendix

% 完工
\end{document}
